\input{preamble}

% OK, start here.
%
\begin{document}

\title{Compléments sur les schémas en groupoïdes}


\maketitle

\phantomsection
\label{section-phantom}

\tableofcontents

\section{Introduction}
\label{section-introduction}

\noindent
Ce chapitre est consacré à des sujets avancés sur les schémas en groupoïdes.
Bien que les résultats soient énoncés en termes de schémas en groupoïdes, le
lecteur doit garder à l'esprit le diagramme $2$-cartésien
\begin{equation}
\label{equation-quotient-stack}
\vcenter{
\xymatrix{
R \ar[r] \ar[d] & U \ar[d] \\
U \ar[r] & [U/R]
}
}
\end{equation}
où $[U/R]$ est le champ quotient, voir
Groupoïdes dans les espaces, Remarque \ref{spaces-groupoids-remark-fundamental-square}.
Nombre de résultats sont motivés par la considération de ce diagramme.
Voir par exemple le très bel article \cite{K-M} de Keel et Mori.





\section{Notations}
\label{section-notation}

\noindent
Nous continuons à suivre les conventions et les notations introduites dans
Groupoïdes, section \ref{groupoids-section-notation}.






\section{Diagrammes utiles}
\label{section-diagrams}

\noindent
Nous rappelons brièvement les résultats des
Lemmes \ref{groupoids-lemma-diagram} et
\ref{groupoids-lemma-diagram-pull} de Groupoïdes,
afin de pouvoir nous y référer aisément dans ce chapitre.
Soit $S$ un schéma.
Soit $(U, R, s, t, c)$ un schéma en groupoïdes sur $S$.
Dans le diagramme commutatif
\begin{equation}
\label{equation-diagram}
\vcenter{
\xymatrix{
& U & \\
R \ar[d]_s \ar[ru]^t &
R \times_{s, U, t} R
\ar[l]^-{\text{pr}_0} \ar[d]^{\text{pr}_1} \ar[r]_-c &
R \ar[d]^s \ar[lu]_t \\
U & R \ar[l]_t \ar[r]^s & U
}
}
\end{equation}
les deux carrés inférieurs sont des carrés de produit fibré.
De plus, le triangle supérieur (qui est en réalité un carré)
est lui aussi cartésien.

\medskip\noindent
Le diagramme
\begin{equation}
\label{equation-pull}
\vcenter{
\xymatrix{
R \times_{t, U, t} R
\ar@<1ex>[r]^-{\text{pr}_1} \ar@<-1ex>[r]_-{\text{pr}_0}
\ar[d]_{\text{pr}_0 \times c \circ (i, 1)} &
R \ar[r]^t \ar[d]^{\text{id}_R} &
U \ar[d]^{\text{id}_U} \\
R \times_{s, U, t} R
\ar@<1ex>[r]^-c \ar@<-1ex>[r]_-{\text{pr}_0} \ar[d]_{\text{pr}_1} &
R \ar[r]^t \ar[d]^s &
U \\
R \ar@<1ex>[r]^s \ar@<-1ex>[r]_t &
U
}
}
\end{equation}
est commutatif. Les deux lignes supérieures sont isomorphes par les applications verticales indiquées.
Les deux carrés inférieurs de gauche sont cartésiens.





\section{Faisceau des différentielles}
\label{section-differentials}

\noindent
Le lemme suivant est l'analogue du
Lemme \ref{groupoids-lemma-group-scheme-module-differentials} de Groupoïdes.

\begin{lemma}
\label{lemma-sheaf-differentials}
Soit $S$ un schéma.
Soit $(U, R, s, t, c)$ un schéma en groupoïdes sur $S$.
Le faisceau des différentielles de $R$, considéré comme un schéma sur
$U$ par $t$, est un quotient de l'image réciproque par $t$ du faisceau conormal de
l'immersion $e : U \to R$. En formule : il existe une surjection canonique
$t^*\mathcal{C}_{U/R} \to \Omega_{R/U}$. Si $s$ est plat, alors
cette application est un isomorphisme.
\end{lemma}

\begin{proof}
Notons que $e : U \to R$ est une immersion, puisqu'il s'agit d'une section
du morphisme $s$, voir
Schémas, Lemme \ref{schemes-lemma-section-immersion}.
Considérons le diagramme suivant
$$
\xymatrix{
R \ar[r]_-{(1, i)} \ar[d]_t &
R \times_{s, U, t} R \ar[d]^c \ar[rr]_{(\text{pr}_0, i \circ \text{pr}_1)} & &
R \times_{t, U, t} R \\
U \ar[r]^e &
R
}
$$
Le carré de gauche est cartésien, car si $a \circ b = e$, alors
$b = i(a)$. La composée des applications horizontales est le
morphisme diagonal de $t : R \to U$. La flèche horizontale supérieure droite est un
isomorphisme. Puisque $\Omega_{R/U}$ est le faisceau conormal de la
composée, il est isomorphe au faisceau conormal de
$(1, i)$. Par le
Lemme \ref{morphisms-lemma-conormal-functorial-flat} de Morphismes,
nous obtenons la surjection $t^*\mathcal{C}_{U/R} \to \Omega_{R/U}$,
et, si $c$ est plat, il s'agit d'un isomorphisme. Puisque $c$ est un changement de base
de $s$ d'après les propriétés du Diagramme (\ref{equation-pull}),
nous concluons que, si $s$ est plat, alors $c$ est plat, voir
Morphismes, Lemme \ref{morphisms-lemma-base-change-flat}.
\end{proof}






\section{Structure locale}
\label{section-local}

\noindent
Soit $S$ un schéma.
Soit $(U, R, s, t, c, e, i)$ un schéma en groupoïdes sur $S$.
Soit $u \in U$ un point. Dans cette section, nous expliquons quelle
structure nous obtenons sur les anneaux locaux
$$
A = \mathcal{O}_{U, u}
\quad\text{et}\quad
B = \mathcal{O}_{R, e(u)}
$$
Nous conviendrons de désigner les homomorphismes d'anneaux locaux
induits par les morphismes $s, t, c, e, i$ par les lettres correspondantes.
En particulier, nous avons un diagramme commutatif
$$
\xymatrix{
A \ar[rd]_t \ar[rrd]^1 \\
& B \ar[r]^e & A \\
A \ar[ru]^s \ar[rru]_1
}
$$
d'anneaux locaux. Ainsi, si $I \subset B$ désigne le noyau de $e : B \to A$,
alors $B = s(A) \oplus I = t(A) \oplus I$. Notons
$$
C = \mathcal{O}_{R \times_{s, U, t} R, (e(u), e(u))}
$$
Nous avons alors
$$
C = (B \otimes_{s, A, t} B)_{\mathfrak m_B \otimes B + B \otimes \mathfrak m_B}
$$
Soit $J \subset C$ l'idéal de $C$ engendré par $I \otimes B + B \otimes I$.
Alors $J$ est également le noyau de l'homomorphisme d'anneaux locaux
$$
(e, e) : C \longrightarrow A
$$
La loi de composition $c : R \times_{s, U, t} R \to R$ correspond à un
homomorphisme d'anneaux
$$
c : B \longrightarrow C
$$
envoyant $I$ dans $J$.

\begin{lemma}
\label{lemma-first-order-structure-c}
L'application $I/I^2 \to J/J^2$ induite par $c$ est la composée
$$
I/I^2 \xrightarrow{(1, 1)} I/I^2 \oplus I/I^2 \to J/J^2
$$
où la seconde flèche provient de l'égalité
$J = (I \otimes B + B \otimes I)C$.
L'application $i : B \to B$ induit l'application $-1 : I/I^2 \to I/I^2$.
\end{lemma}

\begin{proof}
Pour décrire un homomorphisme local de $C$ vers un autre anneau local,
il suffit d'indiquer ce qu'il advient des éléments de la forme
$b_1 \otimes b_2$. En gardant cela à l'esprit, nous avons les deux applications canoniques
$$
e_2 : C \to B,\ b_1 \otimes b_2 \mapsto b_1s(e(b_2)),\quad
e_1 : C \to B,\ b_1 \otimes b_2 \mapsto t(e(b_1))b_2
$$
correspondant aux applications
$R \to R \times_{s, U, t} R$ données par
$r \mapsto (r, e(s(r)))$ et $r \mapsto (e(t(r)), r)$.
Ces applications définissent des applications $J/J^2 \to I/I^2$ qui, ensemble,
donnent un inverse de l'application $I/I^2 \oplus I/I^2 \to J/J^2$
du lemme. Ainsi, pour démontrer l'assertion, il nous suffit de montrer
que $e_1 \circ c : B \to B$ et $e_2 \circ c : B \to B$
sont les applications identiques. Cela résulte du fait que les deux
composées $R \to R \times_{s, U, t} R \to R$ sont des identités.

\medskip\noindent
L'assertion concernant $i$ résulte de celle concernant $c$ et du
fait que $c \circ (1, i) = e \circ t$. Nous omettons quelques détails.
\end{proof}






\section{Propriétés des groupoïdes}
\label{section-technical-lemma}

\noindent
Soit $(U, R, s, t, c)$ un schéma en groupoïdes.
L'idée qui sous-tend les résultats de cette section est que $s: R \to U$
est un changement de base du morphisme $U \to [U/R]$ (voir
le Diagramme (\ref{equation-quotient-stack})).
Les propriétés locales de $s : R \to U$ devraient donc refléter les
propriétés locales du morphisme $U \to [U/R]$.
Cela ne fonctionne pas, car $[U/R]$ n'est pas toujours un champ algébrique et
nous ne pouvons donc pas parler des propriétés géométriques ou algébriques de
$U \to [U/R]$.
Il s'avère toutefois que nous pouvons en récupérer une partie sans même
faire référence au champ quotient.

\medskip\noindent
Voici un premier exemple d'un tel résultat. Grosso modo, l'ouvert $W \subset U'$
trouvé dans le lemme est le lieu où le morphisme
$U' \to [U/R]$ possède la propriété $\mathcal{P}$.

\begin{lemma}
\label{lemma-local-source}
Soit $S$ un schéma.
Soit $(U, R, s, t, c, e, i)$ un groupoïde sur $S$.
Soit $g : U' \to U$ un morphisme de schémas.
Notons $h$ la composée
$$
\xymatrix{
h : U' \times_{g, U, t} R \ar[r]_-{\text{pr}_1} & R \ar[r]_s & U.
}
$$
Soient $\mathcal{P}, \mathcal{Q}, \mathcal{R}$ des propriétés de morphismes
de schémas. Supposons que
\begin{enumerate}
\item $\mathcal{R} \Rightarrow \mathcal{Q}$,
\item $\mathcal{Q}$ est préservée par changement de base et par composition,
\item pour tout morphisme $f : X \to Y$ qui possède $\mathcal{Q}$, il existe un
plus grand ouvert $W(\mathcal{P}, f) \subset X$ tel que $f|_{W(\mathcal{P}, f)}$
possède $\mathcal{P}$, et
\item pour tout morphisme $f : X \to Y$ qui possède $\mathcal{Q}$,
et tout morphisme $Y' \to Y$ qui possède $\mathcal{R}$, nous avons
$Y' \times_Y W(\mathcal{P}, f) = W(\mathcal{P}, f')$, où
$f' : X_{Y'} \to Y'$ est le changement de base de $f$.
\end{enumerate}
Si $s, t$ possèdent $\mathcal{R}$ et si $g$ possède $\mathcal{Q}$, alors
il existe un sous-schéma ouvert $W \subset U'$ tel que
$W \times_{g, U, t} R = W(\mathcal{P}, h)$.
\end{lemma}

\begin{proof}
Notons que le diagramme suivant est commutatif
$$
\xymatrix{
U' \times_{g, U, t} R \times_{t, U, t} R
\ar[rr]_-{\text{pr}_{12}}
\ar@<1ex>[d]^-{\text{pr}_{02}} \ar@<-1ex>[d]_-{\text{pr}_{01}} & &
R \times_{t, U, t} R
\ar@<1ex>[d]^-{\text{pr}_1} \ar@<-1ex>[d]_-{\text{pr}_0}
\\
U' \times_{g, U, t} R \ar[rr]^{\text{pr}_1} & & R
}
$$
dont les deux carrés sont cartésiens (on utilise ici que les deux applications
$t \circ \text{pr}_i : R \times_{t, U, t} R \to U$ sont égales).
En combinant ceci avec les propriétés du diagramme (\ref{equation-pull}),
nous obtenons un diagramme commutatif
$$
\xymatrix{
U' \times_{g, U, t} R \times_{t, U, t} R
\ar[rr]_-{c \circ (i, 1)}
\ar@<1ex>[d]^-{\text{pr}_{02}} \ar@<-1ex>[d]_-{\text{pr}_{01}} & &
R
\ar@<1ex>[d]^-{s} \ar@<-1ex>[d]_-{t}
\\
U' \times_{g, U, t} R \ar[rr]^h & & U
}
$$
dont les deux carrés sont cartésiens.

\medskip\noindent
Supposons que $s, t$ possèdent $\mathcal{R}$ et que $g$ possède $\mathcal{Q}$.
Alors $h$ possède $\mathcal{Q}$, car il est la composée de $s$ (qui possède
$\mathcal{R}$, donc $\mathcal{Q}$) et d'un changement de base de $g$ (qui
possède $\mathcal{Q}$). Ainsi, $W(\mathcal{P}, h) \subset U' \times_{g, U, t} R$
existe. D'après nos hypothèses, nous avons
$\text{pr}_{01}^{-1}(W(\mathcal{P}, h)) =
\text{pr}_{02}^{-1}(W(\mathcal{P}, h))$
car ce sont tous deux le plus grand ouvert sur lequel $c \circ (i, 1)$ possède $\mathcal{P}$.
Notons que la projection $U' \times_{g, U, t} R \to U'$ possède une section, à savoir
$\sigma : U' \to U' \times_{g, U, t} R$, $u' \mapsto (u', e(g(u')))$.
De plus, par l'isomorphisme
$$
(U' \times_{g, U, t} R) \times_{U'} (U' \times_{g, U, t} R)
=
U' \times_{g, U, t} R \times_{t, U, t} R
$$
les deux projections du membre de gauche
vers $U' \times_{g, U, t} R$ coïncident avec les morphismes $\text{pr}_{01}$
et $\text{pr}_{02}$ du membre de droite. Puisque
$\text{pr}_{01}^{-1}(W(\mathcal{P}, h)) =
\text{pr}_{02}^{-1}(W(\mathcal{P}, h))$
nous concluons que $W(\mathcal{P}, h)$ est l'image réciproque d'un sous-ensemble de $U$,
qui est nécessairement l'ouvert
$W = \sigma^{-1}(W(\mathcal{P}, h))$.
\end{proof}

\begin{remark}
\label{remark-local-source-warning}
Avertissement :
le Lemme \ref{lemma-local-source}
doit être utilisé avec précaution.
Il s'applique par exemple à $\mathcal{P}=$``plat'', $\mathcal{Q}=$``sans condition'',
et $\mathcal{R}=$``plat et localement de présentation finie''. Toutefois, pour un
morphisme de schémas $f : X \to Y$, le plus grand ouvert $W \subset X$ tel que
$f|_W$ soit plat n'est {\it pas} l'ensemble des points où $f$ est plat !
\end{remark}

\begin{remark}
\label{remark-local-source-apply}
Malgré l'avertissement de la
Remarque \ref{remark-local-source-warning},
il existe certains cas où le
Lemme \ref{lemma-local-source}
peut être utilisé sans introduire trop d'ambiguïté.
Nous en donnons une liste. Dans chaque cas, nous omettons la vérification des
hypothèses (1) et (2) et donnons des références qui impliquent
(3) et (4). Voici la liste :
\begin{enumerate}
\item $\mathcal{Q} = \mathcal{R} =$``localement de type fini'', et
$\mathcal{P} =$``de dimension relative $\leq d$''.
Voir
Morphismes, Définition \ref{morphisms-definition-relative-dimension-d}
et
Morphismes, Lemmes \ref{morphisms-lemma-openness-bounded-dimension-fibres} et
\ref{morphisms-lemma-dimension-fibre-after-base-change}.
\item $\mathcal{Q} = \mathcal{R} =$``localement de type fini'', et
$\mathcal{P} =$``localement quasi-fini''.
C'est le cas $d = 0$ du point précédent, voir
Morphismes, Lemme \ref{morphisms-lemma-locally-quasi-finite-rel-dimension-0}.
\item $\mathcal{Q} = \mathcal{R} =$``localement de type fini'', et
$\mathcal{P} =$``non ramifié''.
Voir
Morphismes, Lemmes \ref{morphisms-lemma-unramified-characterize} et
\ref{morphisms-lemma-set-points-where-fibres-unramified}.
\end{enumerate}
L'intérêt des cas énumérés ci-dessus est qu'il n'est pas nécessaire de
supposer que $s, t$ sont plats pour obtenir une conclusion sur le lieu
où le morphisme $h$ possède la propriété $\mathcal{P}$. Poursuivons la
liste :
\begin{enumerate}
\item[(4)] $\mathcal{Q} =$``localement de présentation finie'',
$\mathcal{R} =$``plat et localement de présentation finie'', et
$\mathcal{P} =$``plat''. Voir
Compléments sur les morphismes, Théorème
\ref{more-morphisms-theorem-openness-flatness} et
Lemme \ref{more-morphisms-lemma-flat-locus-base-change}.
\item[(5)] $\mathcal{Q} =$``localement de présentation finie'',
$\mathcal{R} =$``plat et localement de présentation finie'', et
$\mathcal{P}=$``de Cohen-Macaulay''. Voir
Compléments sur les morphismes, Définition \ref{more-morphisms-definition-CM}
et
Compléments sur les morphismes, Lemmes \ref{more-morphisms-lemma-base-change-CM} et
\ref{more-morphisms-lemma-flat-finite-presentation-CM-open}.
\item[(6)] $\mathcal{Q} =$``localement de présentation finie'',
$\mathcal{R} =$``plat et localement de présentation finie'', et
$\mathcal{P}=$``syntomique'' ; voir
Morphismes, Lemme \ref{morphisms-lemma-set-points-where-fibres-lci}
(le lieu est automatiquement ouvert).
\item[(7)] $\mathcal{Q} =$``localement de présentation finie'',
$\mathcal{R} =$``plat et localement de présentation finie'', et
$\mathcal{P}=$``lisse''. Voir
Morphismes, Lemme \ref{morphisms-lemma-set-points-where-fibres-smooth}
(le lieu est automatiquement ouvert).
\item[(8)] $\mathcal{Q} =$``localement de présentation finie'',
$\mathcal{R} =$``plat et localement de présentation finie'', et
$\mathcal{P}=$``\'etale''. Voir
Morphismes, Lemme \ref{morphisms-lemma-set-points-where-fibres-etale}
(le lieu est automatiquement ouvert).
\end{enumerate}
\end{remark}

\noindent
Voici le second résultat. Il faut penser à l'ouvert $W \subset U$ invariant par $R$ comme à
l'image réciproque du plus grand ouvert de $[U/R]$ au-dessus duquel
le morphisme $U \to [U/R]$ possède la propriété $\mathcal{P}$.

\begin{lemma}
\label{lemma-property-invariant}
Soit $S$ un schéma.
Soit $(U, R, s, t, c)$ un groupoïde sur $S$.
Soit $\tau \in \{Zariski, \linebreak[0] fppf,
\linebreak[0] \etale, \linebreak[0]
smooth, \linebreak[0] syntomic\}$\footnote{Le fait que $fpqc$ soit absent
n'est pas une coquille.}. Soit $\mathcal{P}$ une propriété de morphismes de schémas
qui est $\tau$-locale sur la base
(Descente, Définition \ref{descent-definition-property-morphisms-local}).
Supposons que $\{s : R \to U\}$ et $\{t : R \to U\}$ soient des recouvrements pour la
topologie $\tau$. Soit $W \subset U$ le sous-schéma ouvert maximal tel que
$s|_{s^{-1}(W)} : s^{-1}(W) \to W$ possède la propriété $\mathcal{P}$.
Alors $W$ est invariant par $R$, voir
Groupoïdes, Définition \ref{groupoids-definition-invariant-open}.
\end{lemma}

\begin{proof}
L'existence et les propriétés de l'ouvert $W \subset U$ sont décrites dans
Descente, Lemme \ref{descent-lemma-largest-open-of-the-base}.
Dans le
Diagramme (\ref{equation-diagram}),
soit $W_1 \subset R$ le sous-schéma ouvert maximal au-dessus duquel le morphisme
$\text{pr}_1 : R \times_{s, U, t} R \to R$ possède la propriété $\mathcal{P}$.
Il résulte du
Lemme \ref{descent-lemma-largest-open-of-the-base} de Descente déjà cité
et de l'hypothèse que $\{s : R \to U\}$ et $\{t : R \to U\}$ sont des recouvrements
pour la topologie $\tau$ que $t^{-1}(W) = W_1 = s^{-1}(W)$, comme souhaité.
\end{proof}

\begin{lemma}
\label{lemma-property-G-invariant}
Soit $S$ un schéma.
Soit $(U, R, s, t, c)$ un groupoïde sur $S$.
Soit $G \to U$ son schéma en groupes stabilisateur.
Soit $\tau \in \{fppf, \linebreak[0] \etale, \linebreak[0]
smooth, \linebreak[0] syntomic\}$.
Soit $\mathcal{P}$ une propriété de morphismes qui est $\tau$-locale
sur la base. Supposons que $\{s : R \to U\}$ et $\{t : R \to U\}$ soient des recouvrements
pour la topologie $\tau$. Soit $W \subset U$ le sous-schéma ouvert maximal
tel que $G_W \to W$ possède la propriété $\mathcal{P}$. Alors $W$ est invariant par $R$
(voir
Groupoïdes, Définition
\ref{groupoids-definition-invariant-open}).
\end{lemma}

\begin{proof}
L'existence et les propriétés de l'ouvert $W \subset U$ sont décrites dans
Descente, Lemme \ref{descent-lemma-largest-open-of-the-base}.
Le morphisme
$$
G \times_{U, t} R \longrightarrow R \times_{s, U} G, \quad
(g, r) \longmapsto (r, r^{-1} \circ g \circ r)
$$
est un isomorphisme au-dessus de $R$ (où $\circ$ désigne
la composition dans le groupoïde). Ainsi, $s^{-1}(W) = t^{-1}(W)$ par les
propriétés de $W$ démontrées dans le
Lemme \ref{descent-lemma-largest-open-of-the-base} de Descente déjà cité.
\end{proof}



\section{Comparaison des fibres}
\label{section-fibres}

\noindent
Soit $(U, R, s, t, c, e, i)$ un schéma en groupoïdes sur $S$.
Le Diagramme (\ref{equation-diagram})
nous fournit un moyen de comparer les fibres de l'application $s : R \to U$ dans un groupoïde.
Pour un point $u \in U$, nous noterons $F_u = s^{-1}(u)$ la fibre
schématique de $s : R \to U$ au-dessus de $u$. Par exemple, le diagramme
implique que, si $u, u' \in U$ sont des points tels
que $s(r) = u$ et $t(r) = u'$, alors
$(F_u)_{\kappa(r)} \cong (F_{u'})_{\kappa(r)}$.
C'est un cas particulier du résultat plus général et plus précis du
Lemme \ref{lemma-two-fibres}
ci-dessous. Pour le voir, prenons $r' = i(r)$.

\medskip\noindent
Un couple $(X, x)$ formé d'un schéma $X$ et d'un point $x \in X$ est parfois
appelé le {\it germe de $X$ en $x$}.
Un {\it morphisme de germes} $f : (X, x) \to (S, s)$
est un morphisme $f : U \to S$ défini sur un voisinage ouvert
de $x$ et tel que $f(x) = s$. Deux tels
$f$, $f'$ sont dits définir le même morphisme de germes
si et seulement si $f$ et $f'$ coïncident sur un voisinage ouvert de $x$.
Soit $\tau \in \{Zariski, \etale, smooth, syntomic, fppf\}$.
Introduisons provisoirement la notion suivante : nous disons que deux morphismes
de germes $f : (X, x) \to (S, s)$ et $f' : (X', x') \to (S', s')$
sont {\it localement isomorphes sur la base pour la topologie $\tau$},
s'il existe un schéma pointé $(S'', s'')$ et des morphismes de germes
$g : (S'', s'') \to (S, s)$ et $g' : (S'', s'') \to (S', s')$
tels que
\begin{enumerate}
\item $g$ et $g'$ sont des immersions ouvertes (resp.\ \'etales, lisses, syntomiques,
plates et localement de présentation finie) en $s''$,
\item il existe un isomorphisme
$$
(S'' \times_{g, S, f} X, \tilde x)
\cong
(S'' \times_{g', S', f'} X', \tilde  x')
$$
de germes au-dessus du germe $(S'', s'')$, pour un certain choix de points
$\tilde x$ et $\tilde x'$ situés au-dessus de $(s'', x)$ et $(s'', x')$.
\end{enumerate}
Enfin, nous disons simplement que les morphismes de germes
$f : (X, x) \to (S, s)$ et $f' : (X', x') \to (S', s')$
sont {\it isomorphes après des changements de base plats} s'il existe
$S'', s'', g, g'$ comme ci-dessus, mais en remplaçant (1) par
la condition que $g$ et $g'$ soient plats en $s''$ (celle-ci est
bien plus faible que toutes les conditions $\tau$ ci-dessus,
car un morphisme plat n'est pas nécessairement ouvert).

\begin{lemma}
\label{lemma-two-fibres}
Soit $S$ un schéma.
Soit $(U, R, s, t, c)$ un groupoïde sur $S$.
Soient $r, r' \in R$ tels que $t(r) = t(r')$ dans $U$.
Posons $u = s(r)$, $u' = s(r')$.
Notons $F_u = s^{-1}(u)$ et $F_{u'} = s^{-1}(u')$ les fibres
schématiques.
\begin{enumerate}
\item Il existe une extension de corps commune
$\kappa(u) \subset k$, $\kappa(u') \subset k$ et
un isomorphisme $(F_u)_k \cong (F_{u'})_k$.
\item On peut choisir l'isomorphisme de (1) de telle sorte qu'un point
situé au-dessus de $r$ s'envoie sur un point situé au-dessus de $r'$.
\item Si les morphismes $s$, $t$ sont plats, alors les morphismes de germes
$s : (R, r) \to (U, u)$ et $s : (R, r') \to (U, u')$ sont
localement isomorphes sur la base pour la topologie plate.
\item Si les morphismes $s$, $t$ sont \'etales
(resp.\ lisses, syntomiques, ou plats et localement de présentation finie),
alors les morphismes de germes $s : (R, r) \to (U, u)$ et
$s : (R, r') \to (U, u')$ sont localement isomorphes sur la base
pour la topologie \'etale (resp.\ lisse, syntomique, ou fppf).
\end{enumerate}
\end{lemma}

\begin{proof}
Nous utilisons à plusieurs reprises les propriétés et l'existence du
diagramme (\ref{equation-diagram}).
D'après les propriétés du diagramme (et
Schémas, lemme \ref{schemes-lemma-points-fibre-product}),
il existe un point $\xi$ de $R \times_{s, U, t} R$
tel que $\text{pr}_0(\xi) = r$ et $c(\xi) = r'$.
Posons $\tilde r = \text{pr}_1(\xi) \in R$.

\medskip\noindent
Démonstration de (1). Posons $k = \kappa(\tilde r)$. Puisque $t(\tilde r) = u$
et $s(\tilde r) = u'$, nous voyons que $k$ est une extension commune
de $\kappa(u)$ et de $\kappa(u')$ et, en fait, que
$(F_u)_k$ et $(F_{u'})_k$ sont tous deux isomorphes à la fibre de
$\text{pr}_1 : R \times_{s, U, t} R \to R$ au-dessus de $\tilde r$.
Ainsi, (1) est démontré.

\medskip\noindent
L'assertion (2) en résulte, puisque le point $\xi$ s'envoie sur $r$, resp.\ $r'$.

\medskip\noindent
L'assertion (3) découle clairement de ce qui précède (en utilisant le point $\xi$ pour
$\tilde u$ et $\tilde u'$) et des définitions.

\medskip\noindent
Si $s$ et $t$ sont plats et de présentation finie, alors
ce sont des morphismes ouverts (Morphismes, lemme \ref{morphisms-lemma-fppf-open}).
L'image d'un certain voisinage ouvert affine $V''$ de $\tilde r$ recouvrira donc
un voisinage ouvert $V$ de $u$, resp.\ $V'$ de $u'$.
On peut les utiliser pour montrer que les propriétés (1) et (2) de la
définition de ``localement isomorphes sur la base pour la
topologie $\tau$''.
\end{proof}







\section{Présentations de Cohen-Macaulay}
\label{section-CM}

\noindent
Étant donné un groupoïde $(U, R, s, t, c)$ avec $s, t$ plats et
localement de présentation finie, il existe un groupoïde ``équivalent''
$(U', R', s', t', c')$ tel que $s'$ et $t'$ soient des morphismes
de Cohen-Macaulay (et localement de présentation finie). Voir
Compléments sur les morphismes, section \ref{more-morphisms-section-CM},
pour plus de renseignements sur les morphismes de Cohen-Macaulay.
Ici, ``équivalent'' peut signifier que les champs quotients
$[U/R]$ et $[U'/R']$ sont des champs équivalents, voir
Groupoïdes dans les espaces, section \ref{spaces-groupoids-section-stacks}
et la section \ref{spaces-groupoids-section-quotient-stack-restrict}.

\begin{lemma}
\label{lemma-make-CM}
Soit $S$ un schéma.
Soit $(U, R, s, t, c)$ un groupoïde au-dessus de $S$.
Supposons $s$ et $t$ plats et localement de présentation finie.
Alors il existe un ouvert $U' \subset U$ tel que
\begin{enumerate}
\item $t^{-1}(U') \subset R$ est le plus grand sous-schéma ouvert de
$R$ sur lequel le morphisme $s$ est de Cohen-Macaulay,
\item $s^{-1}(U') \subset R$ est le plus grand sous-schéma ouvert de
$R$ sur lequel le morphisme $t$ est de Cohen-Macaulay,
\item le morphisme $t|_{s^{-1}(U')} : s^{-1}(U') \to U$ est
surjectif,
\item le morphisme $s|_{t^{-1}(U')} : t^{-1}(U') \to U$ est
surjectif, et
\item la restriction $R' = s^{-1}(U') \cap t^{-1}(U')$
de $R$ à $U'$ définit un groupoïde $(U', R', s', t', c')$ qui a la propriété
que les morphismes $s'$ et $t'$ sont de Cohen-Macaulay et localement de
présentation finie.
\end{enumerate}
\end{lemma}

\begin{proof}
Appliquons le
lemme \ref{lemma-local-source}
avec
$g = \text{id}$ et
$\mathcal{Q} =$``localement de présentation finie'',
$\mathcal{R} =$``plat et localement de présentation finie'', et
$\mathcal{P}=$``de Cohen-Macaulay'', voir
la remarque \ref{remark-local-source-apply}.
Cela nous donne un ouvert $U' \subset U$ tel que
Alors $t^{-1}(U') \subset R$ est le plus grand sous-schéma ouvert de $R$
sur lequel le morphisme $s$ est de Cohen-Macaulay.
Ceci prouve (1).
Soit $i : R \to R$ l'inverse du groupoïde.
Puisque $i$ est un isomorphisme, et $s \circ i = t$ et $t \circ i = s$,
nous voyons que $s^{-1}(U')$ est aussi le plus grand ouvert de $R$ sur lequel $t$ est
de Cohen-Macaulay. Ceci prouve (2).
D'après
Compléments sur les morphismes,
lemme \ref{more-morphisms-lemma-flat-finite-presentation-CM-open},
l'ouvert $t^{-1}(U')$ est dense dans chaque fibre de $s : R \to U$.
Ceci prouve (3). Même argument pour (4).
L'assertion (5) est une conséquence formelle de (1) et (2) et de la discussion
sur les restrictions dans
Groupoïdes, section \ref{groupoids-section-restrict-groupoid}.
\end{proof}








\section{Restriction des groupoïdes}
\label{section-restricting-groupoids}

\noindent
Dans cette section, nous rassemblons plusieurs lemmes sur les
propriétés des groupoïdes qui se transmettent aux restrictions.
La plupart de ces lemmes se démontrent en considérant le
diagramme qui définit
\begin{equation}
\label{equation-restriction}
\vcenter{
\xymatrix{
R' \ar[d] \ar[r] \ar@/_3pc/[dd]_{t'} \ar@/^1pc/[rr]^{s'}&
R \times_{s, U} U' \ar[r] \ar[d] &
U' \ar[d]^g \\
U' \times_{U, t} R \ar[d] \ar[r] &
R \ar[r]^s \ar[d]_t &
U \\
U' \ar[r]^g &
U
}
}
\end{equation}
une restriction. Voir
Groupoïdes, lemme \ref{groupoids-lemma-restrict-groupoid}.

\begin{lemma}
\label{lemma-restrict-preserves-type}
Soit $S$ un schéma.
Soit $(U, R, s, t, c)$ un schéma en groupoïdes au-dessus de $S$.
Soit $g : U' \to U$ un morphisme de schémas.
Soit $(U', R', s', t', c')$ la restriction de
$(U, R, s, t, c)$ par $g$.
\begin{enumerate}
\item Si $s, t$ sont localement de type fini et $g$ est localement de
type fini, alors $s', t'$ sont localement de type fini.
\item Si $s, t$ sont localement de présentation finie et $g$ est localement de
présentation finie, alors $s', t'$ sont localement de présentation finie.
\item Si $s, t$ sont plats et $g$ est plat, alors $s', t'$ sont plats.
\item Ajouter ici.
\end{enumerate}
\end{lemma}

\begin{proof}
La propriété d'être localement de type fini est stable par composition
et par changement de base arbitraire, voir
Morphismes, lemmes \ref{morphisms-lemma-composition-finite-type} et
\ref{morphisms-lemma-base-change-finite-type}.
Par conséquent, (1) résulte clairement du diagramme (\ref{equation-restriction}).
Pour les autres cas, voir
Morphismes, lemmes \ref{morphisms-lemma-composition-finite-presentation},
\ref{morphisms-lemma-base-change-finite-presentation},
\ref{morphisms-lemma-composition-flat}, et
\ref{morphisms-lemma-base-change-flat}.
\end{proof}

\noindent
Le lemme suivant aurait pu servir à démontrer les résultats du lemme
précédent de manière plus uniforme.

\begin{lemma}
\label{lemma-restrict-property}
Soit $S$ un schéma.
Soit $(U, R, s, t, c)$ un schéma en groupoïdes au-dessus de $S$.
Soit $g : U' \to U$ un morphisme de schémas.
Soit $(U', R', s', t', c')$ la restriction de
$(U, R, s, t, c)$ par $g$, et soit
$h = s \circ \text{pr}_1 : U' \times_{g, U, t} R \to U$. Si
$\mathcal{P}$ est une propriété des morphismes de schémas telle que
\begin{enumerate}
\item $h$ possède la propriété $\mathcal{P}$, et
\item $\mathcal{P}$ est préservée par changement de base,
\end{enumerate}
alors $s', t'$ possèdent la propriété $\mathcal{P}$.
\end{lemma}

\begin{proof}
Cela est clair, puisque $s'$ est le changement de base de $h$ d'après le
diagramme (\ref{equation-restriction})
et que $t'$ est isomorphe à $s'$ en tant que morphisme de schémas.
\end{proof}

\begin{lemma}
\label{lemma-double-restrict}
Soit $S$ un schéma.
Soit $(U, R, s, t, c)$ un schéma en groupoïdes au-dessus de $S$.
Soient $g : U' \to U$ et $g' : U'' \to U'$ des morphismes de schémas.
Posons $g'' = g \circ g'$.
Soit $(U', R', s', t', c')$ la restriction de $R$ à $U'$.
Soit $h = s \circ \text{pr}_1 : U' \times_{g, U, t} R \to U$,
soit $h' = s' \circ \text{pr}_1 : U'' \times_{g', U', t} R \to U'$, et
soit $h'' = s \circ \text{pr}_1 : U'' \times_{g'', U, t} R \to U$.
Le diagramme suivant est commutatif
$$
\xymatrix{
U'' \times_{g', U', t} R' \ar[d]^{h'} &
(U' \times_{g, U, t} R) \times_U (U'' \times_{g'', U, t} R)
\ar[l] \ar[r] \ar[d] &
U'' \times_{g'', U, t} R \ar[d]_{h''} \\
U' &
U' \times_{g, U, t} R \ar[l]_{\text{pr}_0} \ar[r]^h &
U
}
$$
ses deux carrés étant cartésiens, et la flèche horizontale supérieure gauche
étant donnée par la règle
$$
\begin{matrix}
(U' \times_{g, U, t} R) \times_U (U'' \times_{g'', U, t} R) &
\longrightarrow &
U'' \times_{g', U', t} R' \\
((u', r_0), (u'', r_1)) &
\longmapsto &
(u'', (c(r_1, i(r_0)), (g'(u''), u')))
\end{matrix}
$$
avec les notations expliquées dans la démonstration.
\end{lemma}

\begin{proof}
Nous allons expliciter ceci en exploitant le point de vue fonctoriel
et en ramenant le lemme à un énoncé sur les flèches dans les restrictions
d'une catégorie groupoïde. Dans la dernière formule du lemme, la
notation $((u', r_0), (u'', r_1))$ désigne un point à valeurs dans $T$ de
$(U' \times_{g, U, t} R) \times_U (U'' \times_{g'', U, t} R)$.
Cela signifie que $u', u'', r_0, r_1$ sont des points à valeurs dans $T$ de $U', U'', R, R$
et que $g(u') = t(r_0)$, $g(g'(u'')) = g''(u'') = t(r_1)$, et
$s(r_0) = s(r_1)$. Il serait plus correct d'écrire ici
$g \circ u' = t \circ r_0$ et ainsi de suite, mais cela rend la notation
encore plus illisible. Si nous considérons $r_1$ et $r_0$ comme des flèches dans
une catégorie groupoïde, nous pouvons représenter ceci par le diagramme
$$
\xymatrix{
t(r_0) = g(u') &
s(r_0) = s(r_1) \ar[l]_{r_0} \ar[r]^-{r_1} &
t(r_1) = g(g'(u''))
}
$$
Ce diagramme montre notamment que la composition
$c(r_1, i(r_0))$ a un sens. Rappelons que
$$
R' = R \times_{(t, s), U \times_S U, g \times g} U' \times_S U'
$$
ainsi, un point à valeurs dans $T$ de $R'$ s'écrit $(r, (u'_0, u'_1))$
avec $t(r) = g(u'_0)$ et $s(r) = g(u'_1)$. En particulier, étant donné
$((u', r_0), (u'', r_1))$ comme ci-dessus, on obtient le point à valeurs dans $T$
$(c(r_1, i(r_0)), (g'(u''), u'))$ de $R'$, car
$t(c(r_1, i(r_0))) = t(r_1) = g(g'(u''))$ et
$s(c(r_1, i(r_0))) = s(i(r_0)) = t(r_0) = g(u')$.
Nous laissons au lecteur le soin de montrer que le carré gauche commute
avec cette définition.

\medskip\noindent
Pour montrer que le carré gauche est cartésien,
supposons donnés $(v'', p')$ et $(v', p)$ qui sont des points à valeurs dans $T$ de
$U'' \times_{g', U', t} R'$ et $U' \times_{g, U, t} R$ avec
$v' = s'(p')$. Cela signifie aussi que $g'(v'') = t'(p')$ et
$g(v') = t(p)$. D'après la discussion ci-dessus, nous savons que nous pouvons écrire
$p' = (r, (u_0', u_1'))$ avec $t(r) = g(u'_0)$ et
$s(r) = g(u'_1)$. Avec cette notation, nous voyons que
$v' = s'(p') = u_1'$ et
$g'(v'') = t'(p') = u_0'$. Voici un diagramme
$$
\xymatrix{
s(p) \ar[r]^-p &
g(v') = g(u'_1) \ar[r]^-r &
g(u'_0) = g(g'(v''))
}
$$
Il s'agit de montrer qu'il existe un unique point à valeurs dans $T$
$((u', r_0), (u'', r_1))$ comme ci-dessus tel que
$v' = u'$, $p = r_0$, $v'' = u''$ et $p' = (c(r_1, i(r_0)), (g'(u''), u'))$.
En comparant les deux diagrammes ci-dessus, il est clair que nous n'avons pas d'autre choix
que de poser
$$
((u', r_0), (u'', r_1)) = ((v', p), (v'', c(r, p))
$$
Quelques détails sont omis.
\end{proof}

\begin{lemma}
\label{lemma-double-restrict-property}
Soit $S$ un schéma.
Soit $(U, R, s, t, c)$ un schéma en groupoïdes au-dessus de $S$.
Soient $g : U' \to U$ et $g' : U'' \to U'$ des morphismes de schémas.
Posons $g'' = g \circ g'$.
Soit $(U', R', s', t', c')$ la restriction de $R$ à $U'$.
Soit $h = s \circ \text{pr}_1 : U' \times_{g, U, t} R \to U$,
soit $h' = s' \circ \text{pr}_1 : U'' \times_{g', U', t} R \to U'$, et
soit $h'' = s \circ \text{pr}_1 : U'' \times_{g'', U, t} R \to U$.
Soit $\tau \in \{Zariski, \linebreak[0] \etale, \linebreak[0]
smooth, \linebreak[0] syntomic, \linebreak[0] fppf, \linebreak[0] fpqc\}$. Soit
$\mathcal{P}$ une propriété des morphismes de schémas
qui est préservée par changement de base et qui
est locale sur la base pour la topologie $\tau$. Si
\begin{enumerate}
\item $h(U' \times_U R)$ est ouvert dans $U$,
\item $\{h : U' \times_U R \to h(U' \times_U R)\}$ est un $\tau$-recouvrement,
\item $h'$ possède la propriété $\mathcal{P}$,
\end{enumerate}
alors $h''$ possède la propriété $\mathcal{P}$. Réciproquement, si
\begin{enumerate}
\item[(a)] $\{t : R \to U\}$ est un $\tau$-recouvrement,
\item[(d)] $h''$ possède la propriété $\mathcal{P}$,
\end{enumerate}
alors $h'$ possède la propriété $\mathcal{P}$.
\end{lemma}

\begin{proof}
Cela résulte formellement des propriétés du diagramme du
lemme \ref{lemma-double-restrict}.
Dans le premier cas, remarquons que l'image du morphisme
$h''$ est contenue dans l'image de $h$, puisque $g'' = g \circ g'$.
On peut donc remplacer le $U$ situé dans le coin inférieur droit du
diagramme par $h(U' \times_U R)$. Cela explique l'importance des
conditions (1) et (2) du lemme. Dans le second cas, remarquons que
$\{\text{pr}_0 : U' \times_{g, U, t} R \to U'\}$ est un $\tau$-recouvrement
en tant que changement de base de $\tau$ et de la condition (a).
\end{proof}






\section{Propriétés des groupoïdes sur les corps}
\label{section-properties-groupoids-on-fields}

\noindent
Un ``groupoïde sur un corps'' désigne un schéma en groupoïdes $(U, R, s, t, c)$
où $U$ est le spectre d'un corps. Cela ne signifie {\bf pas} que
$(U, R, s, t, c)$ est défini sur un corps; plus précisément, cela ne
signifie {\bf pas} que les morphismes $s, t : R \to U$ sont égaux.
Étant donné un corps $k$, un groupe abstrait $G$ et un homomorphisme de groupes
$\varphi : G \to \text{Aut}(k)$, on obtient un schéma en groupoïdes
$(U, R, s, t, c)$ au-dessus de $\mathbf{Z}$ en posant
\begin{align*}
U & = \Spec(k) \\
R & = \coprod\nolimits_{g \in G} \Spec(k) \\
s & = \coprod\nolimits_{g \in G} \Spec(\text{id}_k) \\
t & = \coprod\nolimits_{g \in G} \Spec(\varphi(g)) \\
c & = \text{composition dans }G
\end{align*}
Cet exemple est encore un schéma en groupoïdes au-dessus de $\Spec(k^G)$.
Par conséquent, si $G$ est fini, alors $U = \Spec(k)$ est fini sur
$\Spec(k^G)$.
En un certain sens, notre but dans cette section est de montrer que des
conditions de finitude convenables sur $s, t$ imposent à tout groupoïde sur un corps
d'être défini sur un sous-corps d'indice fini $k' \subset k$.

\medskip\noindent
Si $k$ est un corps et $(G, m)$ un schéma en groupes sur $k$, de morphisme
structural $p : G \to \Spec(k)$, alors $(\Spec(k), G, p, p, m)$
est un exemple de groupoïde sur un corps (et, dans ce cas, bien entendu, toute la
structure est définie sur un corps). Ainsi, cette section peut être considérée comme
l'analogue de
Groupoïdes, section \ref{groupoids-section-properties-group-schemes-field}.

\begin{lemma}
\label{lemma-groupoid-on-field-open-multiplication}
Soit $S$ un schéma. Soit $(U, R, s, t, c)$ un schéma en groupoïdes
au-dessus de $S$. Si $U$ est le spectre d'un corps, alors le morphisme de composition
$c : R \times_{s, U, t} R \to R$ est ouvert.
\end{lemma}

\begin{proof}
La composition est isomorphe à la projection
$\text{pr}_1 : R \times_{t, U, t} R \to R$ d'après le
diagramme (\ref{equation-pull}).
La projection est ouverte d'après
Morphismes, lemme \ref{morphisms-lemma-scheme-over-field-universally-open}.
\end{proof}

\begin{lemma}
\label{lemma-groupoid-on-field-separated}
Soit $S$ un schéma. Soit $(U, R, s, t, c)$ un schéma en groupoïdes
au-dessus de $S$. Si $U$ est le spectre d'un corps,
alors $R$ est un schéma séparé.
\end{lemma}

\begin{proof}
D'après
Groupoïdes, lemme \ref{groupoids-lemma-group-scheme-over-field-separated},
le schéma en groupes stabilisateur $G \to U$ est séparé. D'après
Groupoïdes, lemme \ref{groupoids-lemma-diagonal},
le morphisme $j = (t, s) : R \to U \times_S U$ est séparé.
Comme $U$ est le spectre d'un corps, le schéma
$U \times_S U$ est affine (par la construction des produits fibrés dans
Schémas, section \ref{schemes-section-fibre-products}).
Par conséquent, $R$ est un schéma séparé, voir
Schémas, lemme \ref{schemes-lemma-separated-permanence}.
\end{proof}

\begin{lemma}
\label{lemma-groupoid-on-field-homogeneous}
Soit $S$ un schéma. Soit $(U, R, s, t, c)$ un schéma en groupoïdes
au-dessus de $S$. Supposons $U = \Spec(k)$, où $k$ est un corps.
Pour tous points $r, r' \in R$, il existe une extension de corps
$k'/k$ et des points
$r_1, r_2 \in R \times_{s, \Spec(k)} \Spec(k')$
et un diagramme
$$
\xymatrix{
R &
R \times_{s, \Spec(k)} \Spec(k')
\ar[l]_-{\text{pr}_0} \ar[r]^\varphi &
R \times_{s, \Spec(k)} \Spec(k')
\ar[r]^-{\text{pr}_0} &
R
}
$$
tel que $\varphi$ soit un isomorphisme de schémas au-dessus de $\Spec(k')$, et que
l'on ait $\varphi(r_1) = r_2$, $\text{pr}_0(r_1) = r$, et
$\text{pr}_0(r_2) = r'$.
\end{lemma}

\begin{proof}
C'est un cas particulier du
lemme \ref{lemma-two-fibres},
parties (1) et (2).
\end{proof}

\begin{lemma}
\label{lemma-restrict-groupoid-on-field}
Soit $S$ un schéma. Soit $(U, R, s, t, c)$ un schéma en groupoïdes
au-dessus de $S$. Supposons $U = \Spec(k)$, où $k$ est un corps.
Soit $k'/k$ une extension de corps, $U' = \Spec(k')$,
et soit $(U', R', s', t', c')$ la restriction de
$(U, R, s, t, c)$ par $U' \to U$. Dans le diagramme qui la définit
$$
\xymatrix{
R' \ar[d] \ar[r] \ar@/_3pc/[dd]_{t'} \ar@/^1pc/[rr]^{s'} \ar@{..>}[rd] &
R \times_{s, U} U' \ar[r] \ar[d] &
U' \ar[d] \\
U' \times_{U, t} R \ar[d] \ar[r] &
R \ar[r]^s \ar[d]_t &
U \\
U' \ar[r] &
U
}
$$
tous les morphismes sont surjectifs, plats et universellement ouverts.
La flèche pointillée $R' \to R$ est en outre affine.
\end{lemma}

\begin{proof}
Le morphisme $U' \to U$ est égal à $\Spec(k') \to \Spec(k)$;
il est donc affine, surjectif et plat. Les morphismes $s, t : R \to U$
ainsi que le morphisme $U' \to U$ sont universellement ouverts d'après
Morphismes, lemme \ref{morphisms-lemma-scheme-over-field-universally-open}.
Comme $R$ n'est pas vide et que $U$ est le spectre d'un corps, les morphismes
$s, t : R \to U$ sont surjectifs et plats. On conclut alors à l'aide de
Morphismes, lemmes \ref{morphisms-lemma-base-change-surjective},
\ref{morphisms-lemma-composition-surjective},
\ref{morphisms-lemma-composition-open},
\ref{morphisms-lemma-base-change-affine},
\ref{morphisms-lemma-composition-affine},
\ref{morphisms-lemma-base-change-flat}, et
\ref{morphisms-lemma-composition-flat}.
\end{proof}

\begin{lemma}
\label{lemma-groupoid-on-field-explain-points}
Soit $S$ un schéma. Soit $(U, R, s, t, c)$ un schéma en groupoïdes
au-dessus de $S$. Supposons $U = \Spec(k)$, où $k$ est un corps.
Pour tout point $r \in R$, il existe
\begin{enumerate}
\item une extension de corps $k'/k$ avec $k'$ algébriquement clos,
\item un point $r' \in R'$ où $(U', R', s', t', c')$ est la
restriction de $(U, R, s, t, c)$ par $\Spec(k') \to \Spec(k)$
\end{enumerate}
de telle sorte que
\begin{enumerate}
\item le point $r'$ a pour image $r$ par le morphisme $R' \to R$, et
\item les applications $s', t' : R' \to \Spec(k')$ induisent des isomorphismes
$k' \to \kappa(r')$.
\end{enumerate}
\end{lemma}

\begin{proof}
En traduisant l'énoncé géométrique en un énoncé sur les corps,
cela signifie que nous pouvons trouver un diagramme
$$
\xymatrix{
k' & k' \ar[l]^1 & \\
k' \ar[u]^\tau & \kappa(r) \ar[lu]^\sigma & k \ar[l]^-s \ar[lu]_i \\
& k \ar[lu]^i \ar[u]_t
}
$$
où $i : k \to k'$ est le plongement de $k$ dans $k'$,
les applications $s, t : k \to \kappa(r)$ sont induites par $s, t : R \to U$, et
l'application $\tau : k' \to k'$ est un automorphisme. Pour construire un tel
diagramme, on peut procéder comme suit :
\begin{enumerate}
\item Choisir un homomorphisme de corps $i : k \to k'$, avec $k'$ algébriquement clos et de
très grand degré de transcendance sur $k$.
\item Choisir un plongement $\sigma : \kappa(r) \to k'$ tel que
$\sigma \circ s = i$. Un tel $\sigma$ existe, car il suffit de
choisir une base de transcendance $\{x_\alpha\}_{\alpha \in A}$ de $\kappa(r)$
sur $k$ et des $y_\alpha \in k'$, $\alpha \in A$, qui soient algébriquement
indépendants sur $i(k)$, puis d'envoyer $s(k)(\{x_\alpha\})$ dans $k'$
selon les règles $s(\lambda) \mapsto i(\lambda)$ pour $\lambda \in k$
et $x_\alpha \mapsto y_\alpha$ pour $\alpha \in A$.
Prolonger ensuite en $\tau : \kappa(\alpha) \to k'$ en utilisant le fait que $k'$ est
algébriquement clos.
\item Choisir un automorphisme $\tau : k' \to k'$ tel que
$\tau \circ i = \sigma \circ t$. Pour cela, choisir une base de transcendance
$\{x_\alpha\}_{\alpha \in A}$ de $k$ sur son sous-corps premier.
D'une part, compléter $\{i(x_\alpha)\}$ en une base de transcendance de
$k'$ en ajoutant $\{y_\beta\}_{\beta \in B}$, et compléter
$\{\sigma(t(x_\alpha))\}$ en une base de transcendance de $k'$ en ajoutant
$\{z_\gamma\}_{\gamma \in C}$.
Comme $k'$ est algébriquement clos, nous pouvons prolonger l'isomorphisme
$\sigma \circ t \circ i^{-1} : i(k) \to \sigma(t(k))$
en un isomorphisme $\tau' : \overline{i(k)} \to \overline{\sigma(t(k))}$
entre leurs clôtures algébriques dans $k'$.
Comme $k'$ a un grand degré de transcendance,
nous voyons que les ensembles $B$ et $C$ ont le même cardinal.
Ainsi, nous pouvons utiliser une bijection
$B \to C$ pour prolonger $\tau'$ en un isomorphisme
$$
\overline{i(k)}(\{y_\beta\})
\longrightarrow
\overline{\sigma(t(k))}(\{z_\gamma\})
$$
et alors, puisque $k'$ est la clôture algébrique des deux membres, nous
voyons que cela se prolonge en un automorphisme $\tau : k' \to k'$
comme souhaité.
\end{enumerate}
Cela démontre le lemme.
\end{proof}

\begin{lemma}
\label{lemma-groupoid-on-field-move-point}
Soit $S$ un schéma. Soit $(U, R, s, t, c)$ un schéma en groupoïdes
sur $S$. Supposons $U = \Spec(k)$, où $k$ est un corps.
Si $r \in R$ est un point tel que $s, t$ induisent des
isomorphismes $k \to \kappa(r)$, alors l'application
$$
R \longrightarrow R, \quad
x \longmapsto c(r, x)
$$
(voir la démonstration pour une notation précise) est un automorphisme $R \to R$
qui envoie $e$ sur $r$.
\end{lemma}

\begin{proof}
Cela est tout à fait évident si l'on envisage les groupoïdes de manière
fonctorielle. Mais nous allons aussi l'expliciter entièrement.
Notons $a : U \to R$ le morphisme d'image $r$ tel que
$s \circ a = \text{id}_U$, qui existe par l'hypothèse
selon laquelle $s : k \to \kappa(r)$ est un isomorphisme. De même, notons
$b : U \to R$ le morphisme d'image $r$ tel que
$t \circ b = \text{id}_U$. Remarquons que
$b = a \circ (t \circ a)^{-1}$, et, en particulier,
$a \circ s \circ b = b$.

\medskip\noindent
Considérons le morphisme $\Psi : R \to R$ défini sur les points à valeurs dans $T$
par
$$
(f : T \to R) \longmapsto (c(a \circ t \circ f, f) : T \to R)
$$
Pour voir qu'il est défini, nous devons vérifier que
$s \circ a \circ t \circ f = t \circ f$, ce qui est évident puisque $s \circ a = 1$.
Remarquons que $\Phi(e) = a$ ; ainsi, pour démontrer le lemme, il
suffit de montrer que $\Phi$ est un automorphisme de $R$.
Soit $\Phi : R \to R$ le morphisme défini sur les points à valeurs dans $T$ par
$$
(g : T \to R) \longmapsto (c(i \circ b \circ t \circ g, g) : T \to R).
$$
Il est défini car
$s \circ i \circ b \circ t \circ g = t \circ b \circ t \circ g =
t \circ g$. Nous affirmons que $\Phi$ et $\Psi$ sont inverses
l'un de l'autre. Pour le voir, calculons
\begin{align*}
& c(a \circ t \circ c(i \circ b \circ t \circ g, g),
c(i \circ b \circ t \circ g, g)) \\
& =
c(a \circ t \circ i \circ b \circ t \circ g,
c(i \circ b \circ t \circ g, g)) \\
& =
c(a \circ s \circ b \circ t \circ g,
c(i \circ b \circ t \circ g, g)) \\
& =
c(b \circ t \circ g, c(i \circ b \circ t \circ g, g)) \\
& =
c(c(b \circ t \circ g, i \circ b \circ t \circ g), g)) \\
& =
c(e, g) \\
& = g
\end{align*}
où nous avons utilisé la relation $a \circ s \circ b = b$ établie ci-dessus.
Dans l'autre sens, nous avons
\begin{align*}
& c(i \circ b \circ t \circ c(a \circ t \circ f, f), c(a \circ t \circ f, f)) \\
& =
c(i \circ b \circ t \circ a \circ t \circ f, c(a \circ t \circ f, f)) \\
& =
c(i \circ a \circ (t \circ a)^{-1} \circ t \circ a \circ t \circ f,
c(a \circ t \circ f, f)) \\
& =
c(i \circ a \circ t \circ f, c(a \circ t \circ f, f)) \\
& =
c(c(i \circ a \circ t \circ f, a \circ t \circ f), f) \\
& =
c(e, f) \\
& = f
\end{align*}
Le lemme est démontré.
\end{proof}

\begin{lemma}
\label{lemma-groupoid-on-field-translate-open}
Soit $S$ un schéma. Soit $(U, R, s, t, c)$ un schéma en groupoïdes
sur $S$. Si $U$ est le spectre d'un corps, si $W \subset R$ est ouvert
et si $Z \to R$ est un morphisme de schémas, alors l'image de la
composée $Z \times_{s, U, t} W \to R \times_{s, U, t} R \to R$ est ouverte.
\end{lemma}

\begin{proof}
Écrivons $U = \Spec(k)$. Considérons une extension de corps $k'/k$. Notons
$U' = \Spec(k')$. Soit $R'$ la restriction de $R$ par $U' \to U$.
Posons $Z' = Z \times_R R'$ et $W' = R' \times_R W$.
Considérons un point $\xi = (z, w)$ de $Z \times_{s, U, t} W$.
Soit $r \in R$ l'image de $z$ par $Z \to R$.
Choisissons $k' \supset k$ et $r' \in R'$ comme dans le
lemme \ref{lemma-groupoid-on-field-explain-points}.
Nous pouvons choisir $z' \in Z'$ se projetant sur $z$ et sur $r'$.
Nous pouvons alors trouver $\xi' \in Z' \times_{s', U', t'} W'$
se projetant sur $z'$ et sur $\xi$. L'ouvert $c(r', W')$
(lemme \ref{lemma-groupoid-on-field-move-point}) est
contenu dans l'image de $Z' \times_{s', U', t'} W' \to R'$.
Observons que $Z' \times_{s', U', t'} W' = (Z \times_{s, U, t} W)
\times_{R \times_{s, U, t} R} (R' \times_{s', U', t'} R')$.
Par conséquent, l'image de $Z' \times_{s', U', t'} W' \to R' \to R$
est contenue dans l'image de $Z \times_{s, U, t} W \to R$.
Comme le morphisme $R' \to R$ est ouvert (lemme \ref{lemma-restrict-groupoid-on-field}),
nous concluons que l'image contient un voisinage ouvert de
l'image de $\xi$, comme souhaité.
\end{proof}

\begin{lemma}
\label{lemma-groupoid-on-field-geometrically-irreducible}
Soit $S$ un schéma. Soit $(U, R, s, t, c)$ un schéma en groupoïdes
sur $S$. Supposons $U = \Spec(k)$, où $k$ est un corps.
Par abus de notation, notons $e \in R$ l'image du morphisme identité
$e : U \to R$. Alors :
\begin{enumerate}
\item tout anneau local $\mathcal{O}_{R, r}$ de $R$ possède un unique
idéal premier minimal ;
\item il existe exactement une composante irréductible $Z$ de $R$
passant par $e$ ;
\item $Z$ est géométriquement irréductible sur $k$ pour l'un ou l'autre
des morphismes $s$ et $t$.
\end{enumerate}
\end{lemma}

\begin{proof}
Soit $r \in R$ un point.
Dans cette démonstration, nous utiliserons sans autre mention la correspondance entre les composantes irréductibles
de $R$ passant par un point $r$ et les idéaux premiers minimaux de l'anneau
local $\mathcal{O}_{R, r}$.
Choisissons $k \subset k'$ et $r' \in R'$ comme dans le
lemme \ref{lemma-groupoid-on-field-explain-points}.
Remarquons que $\mathcal{O}_{R, r} \to \mathcal{O}_{R', r'}$
est fidèlement plat et local, voir le
lemme \ref{lemma-restrict-groupoid-on-field}.
Ainsi, le résultat pour $r' \in R'$ entraîne le résultat pour $r \in R$.
Autrement dit, nous pouvons supposer que $s, t : k \to \kappa(r)$
sont des isomorphismes. D'après le
lemme \ref{lemma-groupoid-on-field-move-point},
il existe un automorphisme envoyant $e$ sur $r$.
Nous pouvons donc supposer $r = e$, c'est-à-dire que (1) découle de (2).

\medskip\noindent
Démontrons d'abord (2) dans le cas où $k$ est séparablement clos.
Soient en effet $X, Y \subset R$ des composantes irréductibles
passant par $e$. Alors, d'après
Variétés, lemme \ref{varieties-lemma-bijection-irreducible-components} et
\ref{varieties-lemma-separably-closed-irreducible},
le schéma $X \times_{s, U, t} Y$ est lui aussi irréductible.
Ainsi, $c(X \times_{s, U, t} Y) \subset R$ est un sous-ensemble irréductible.
Nous affirmons qu'il contient à la fois $X$ et $Y$ (comme sous-ensembles de $R$).
En effet, soit $T$ le spectre d'un corps. Si $x : T \to X$ est un $T$-point
de $X$, alors $c(x, e \circ s \circ x) = x$ et $e \circ s \circ x$
se factorise par $Y$ puisque $e \in Y$. Il en va de même pour les points de $Y$.
Cela implique manifestement que $X = Y$, c'est-à-dire qu'il existe une unique composante irréductible
de $R$ passant par $e$.

\medskip\noindent
Démonstration de (2) et (3) dans le cas général.
Soit $k \subset k'$ une clôture séparable, et
soit $(U', R', s', t', c')$ la restriction de
$(U, R, s, t, c)$ par $\Spec(k') \to \Spec(k)$.
D'après le paragraphe précédent, il existe exactement une composante irréductible
$Z'$ de $R'$ passant par $e'$.
Notons $e'' \in R \times_{s, U} U'$ le changement de base de $e$.
Comme $R' \to R \times_{s, U} U'$ est fidèlement plat, voir le
lemme \ref{lemma-restrict-groupoid-on-field},
et que $e' \mapsto e''$, il existe exactement une composante
irréductible $Z''$ de $R \times_{s, k} k'$ passant
par $e''$. Comme $R \times_k k' \to R$ est fidèlement
plat, cela implique qu'il existe exactement une composante irréductible $Z$ de $R$
passant par $e$. Cela démontre (2).

\medskip\noindent
Pour démontrer (3), soit $Z''' \subset R \times_k k'$ une composante irréductible quelconque
de $Z \times_k k'$. D'après
Variétés, lemme \ref{varieties-lemma-orbit-irreducible-components},
nous avons $Z''' = \sigma(Z'')$ pour un certain $\sigma \in \text{Gal}(k'/k)$.
Comme $\sigma(e'') = e''$, nous avons $e'' \in Z'''$, et donc
$Z''' = Z''$. Cela signifie que $Z$ est géométriquement irréductible
sur $\Spec(k)$ pour le morphisme $s$.
Le même argument montre que $Z$ est géométriquement irréductible
sur $\Spec(k)$ pour le morphisme $t$.
\end{proof}

\begin{lemma}
\label{lemma-groupoid-on-field-locally-finite-type-dimension}
Soit $S$ un schéma. Soit $(U, R, s, t, c)$ un schéma en groupoïdes
sur $S$. Supposons $U = \Spec(k)$, où $k$ est un corps.
Supposons que $s, t$ soient localement de type fini.
Alors :
\begin{enumerate}
\item $R$ est équidimensionnel ;
\item $\dim(R) = \dim_r(R)$ pour tout $r \in R$ ;
\item pour tout $r \in R$, nous avons
$\text{trdeg}_{s(k)}(\kappa(r)) = \text{trdeg}_{t(k)}(\kappa(r))$ ; et
\item pour tout point fermé $r \in R$, nous avons
$\dim(R) = \dim(\mathcal{O}_{R, r})$.
\end{enumerate}
\end{lemma}

\begin{proof}
Soient $r, r' \in R$.
Alors $\dim_r(R) = \dim_{r'}(R)$ d'après le
lemme \ref{lemma-groupoid-on-field-homogeneous}
et
Morphismes, lemme \ref{morphisms-lemma-dimension-fibre-after-base-change}.
D'après
Morphismes, lemme \ref{morphisms-lemma-dimension-fibre-at-a-point},
nous avons
$$
\dim_r(R) =
\dim(\mathcal{O}_{R, r}) + \text{trdeg}_{s(k)}(\kappa(r)) =
\dim(\mathcal{O}_{R, r}) + \text{trdeg}_{t(k)}(\kappa(r)).
$$
D'autre part, la dimension de $R$ (ou de tout ouvert de $R$)
est la borne supérieure des dimensions des anneaux locaux de $R$, voir
Propriétés, lemme \ref{properties-lemma-codimension-local-ring}.
La dimension de l'anneau local est manifestement maximale aux points fermés $r$, auquel cas
$\text{trdeg}_k(\kappa(r)) = 0$ (par le théorème des zéros de Hilbert, voir
Morphismes, section \ref{morphisms-section-points-finite-type}).
Le lemme en résulte.
\end{proof}

\begin{lemma}
\label{lemma-groupoid-on-field-dimension-equal-stabilizer}
Soit $S$ un schéma. Soit $(U, R, s, t, c)$ un schéma en groupoïdes
sur $S$. Supposons $U = \Spec(k)$, où $k$ est un corps.
Supposons que $s, t$ soient localement de type fini.
Alors $\dim(R) = \dim(G)$, où $G$ est le schéma en groupes stabilisateur de $R$.
\end{lemma}

\begin{proof}
Soit $Z \subset R$ la composante irréductible passant par $e$ (voir le
lemme \ref{lemma-groupoid-on-field-geometrically-irreducible}),
considérée comme un sous-schéma fermé intègre de $R$.
Soit $k'_s$, resp.\ $k'_t$, la clôture intégrale de
$s(k)$, resp.\ $t(k)$, dans $\Gamma(Z, \mathcal{O}_Z)$.
Rappelons que $k'_s$ et $k'_t$ sont des corps, voir
Variétés, lemme \ref{varieties-lemma-integral-closure-ground-field}.
D'après
Variétés, proposition \ref{varieties-proposition-unique-base-field},
nous avons $k'_s = k'_t$ comme sous-anneaux de $\Gamma(Z, \mathcal{O}_Z)$.
Comme $e$ se factorise par $Z$, nous obtenons un diagramme commutatif
$$
\xymatrix{
k \ar[rd]_t \ar[rrd]^1 \\
& \Gamma(Z, \mathcal{O}_Z) \ar[r]^e & k \\
k \ar[ru]^s \ar[rru]_1
}
$$
Cela montre d'une part que $k'_s = s(k)$ et $k'_t = t(k)$, donc que
$s(k) = t(k)$, ce qui, combiné au diagramme ci-dessus, implique
que $s = t$ ! Autrement dit, nous concluons que $Z$ est un sous-schéma fermé
de $G = R \times_{(t, s), U \times_S U, \Delta} U$.
Le lemme en résulte, car $G$ et $R$ sont tous deux équidimensionnels, voir
lemme \ref{lemma-groupoid-on-field-locally-finite-type-dimension} et
Groupoïdes, lemme \ref{groupoids-lemma-group-scheme-finite-type-field}.
\end{proof}

\begin{remark}
\label{remark-warn-dimension-groupoid-on-field}
Avertissement :
Le lemme \ref{lemma-groupoid-on-field-dimension-equal-stabilizer}
est faux sans la condition que $s$ et $t$ soient localement de
type fini.
Un exemple simple consiste à partir de l'action
$$
\mathbf{G}_{m, \mathbf{Q}} \times_{\mathbf{Q}} \mathbf{A}^1_{\mathbf{Q}}
\to \mathbf{A}^1_{\mathbf{Q}}
$$
et à restreindre le schéma en groupoïdes correspondant au point générique de
$\mathbf{A}^1_{\mathbf{Q}}$. Autrement dit, effectuons la restriction suivant le morphisme
$\Spec(\mathbf{Q}(x)) \to
\Spec(\mathbf{Q}[x]) = \mathbf{A}^1_{\mathbf{Q}}$.
On obtient alors un schéma en groupoïdes
$(U, R, s, t, c)$ avec
$U = \Spec(\mathbf{Q}(x))$
et
$$
R = \Spec\left(
\mathbf{Q}(x)[y]\left[
\frac{1}{P(xy)}, P \in \mathbf{Q}[T], P \not = 0
\right]
\right)
$$
Dans ce cas, $\dim(R) = 1$ et $\dim(G) = 0$.
\end{remark}

\begin{lemma}
\label{lemma-groupoid-characteristic-zero-smooth}
Soit $S$ un schéma. Soit $(U, R, s, t, c)$ un schéma en groupoïdes
sur $S$. Supposons que les conditions suivantes soient satisfaites :
\begin{enumerate}
\item $U = \Spec(k)$, où $k$ est un corps ;
\item $s, t$ sont localement de type fini ; et
\item la caractéristique de $k$ est nulle.
\end{enumerate}
Alors $s, t : R \to U$ sont lisses.
\end{lemma}

\begin{proof}
D'après le
lemme \ref{lemma-sheaf-differentials},
le faisceau des différentielles de $R \to U$ est libre.
La lissité résulte donc de
Variétés, lemme \ref{varieties-lemma-char-zero-differentials-free-smooth}.
\end{proof}

\begin{lemma}
\label{lemma-reduced-group-scheme-perfect-field-characteristic-p-smooth}
Soit $S$ un schéma. Soit $(U, R, s, t, c)$ un schéma en groupoïdes
sur $S$. Supposons que les conditions suivantes soient satisfaites :
\begin{enumerate}
\item $U = \Spec(k)$, où $k$ est un corps ;
\item $s, t$ sont localement de type fini ;
\item $R$ est réduit ; et
\item $k$ est parfait.
\end{enumerate}
Alors $s, t : R \to U$ sont lisses.
\end{lemma}

\begin{proof}
D'après le
lemme \ref{lemma-sheaf-differentials},
le faisceau $\Omega_{R/U}$ est libre. Le lemme résulte donc de
Variétés, lemme \ref{varieties-lemma-char-p-differentials-free-smooth}.
\end{proof}










\section{Morphismes de groupoïdes sur des corps}
\label{section-morphisms-groupoids-on-fields}

\noindent
Cette section étudie les morphismes entre groupoïdes sur des corps.
C'est légèrement plus général, mais très proche de l'étude des
morphismes de schémas en groupes sur un corps.

\begin{situation}
\label{situation-morphism-groupoids-on-field}
Soit $S$ un schéma.
Soit $U = \Spec(k)$ un schéma sur $S$, où $k$ est un corps.
Soient $(U, R_1, s_1, t_1, c_1)$ et $(U, R_2, s_2, t_2, c_2)$ des schémas en groupoïdes
sur $S$ ayant la même première composante. Soit $a : R_1 \to R_2$ un morphisme
tel que $(\text{id}_U, a)$ définisse un morphisme de schémas en
groupoïdes sur $S$, voir
Groupoïdes, définition \ref{groupoids-definition-groupoid}.
En particulier, les diagrammes suivants sont commutatifs
$$
\vcenter{
\xymatrix{
R_1 \ar[rrd]^{t_1} \ar[rdd]_{s_1} \ar[rd]_a \\
& R_2 \ar[d]^{t_2} \ar[r]_{s_2} & U \\
& U
}
}
\quad\quad
\vcenter{
\xymatrix{
R_1 \times_{s_1, U, t_1} R_1 \ar[r]_-{c_1} \ar[d]_{a \times a} &
R_1 \ar[d]^a \\
R_2 \times_{s_2, U, t_2} R_2 \ar[r]^-{c_2} &
R_2
}
}
$$
\end{situation}

\noindent
Le lemme suivant généralise
Groupoïdes, lemme \ref{groupoids-lemma-open-subgroup-closed-over-field}.

\begin{lemma}
\label{lemma-open-image-is-closed}
Notations et hypothèses comme dans la
situation \ref{situation-morphism-groupoids-on-field}.
Si $a(R_1)$ est ouvert dans $R_2$, alors $a(R_1)$ est fermé dans $R_2$.
\end{lemma}

\begin{proof}
Soit $r_2 \in R_2$ un point appartenant à l'adhérence de $a(R_1)$.
Nous voulons montrer que $r_2 \in a(R_1)$. Choisissons $k \subset k'$ et
$r_2' \in R'_2$ adapté à $(U, R_2, s_2, t_2, c_2)$ et à $r_2$ comme dans le
lemme \ref{lemma-groupoid-on-field-explain-points}.
Soit $R_i'$ la restriction de $R_i$ par le morphisme
$U' = \Spec(k') \to U = \Spec(k)$.
Soit $a' : R'_1 \to R_2'$ le changement de base de $a$. Le diagramme
$$
\xymatrix{
R'_1 \ar[r]_{a'} \ar[d]_{p_1} & R'_2 \ar[d]^{p_2} \\
R_1 \ar[r]^a & R_2
}
$$
est un carré cartésien. L'image de $a'$ est donc l'image réciproque de
l'image de $a$ par le morphisme $p_2 : R'_2 \to R_2$. D'après le
lemme \ref{lemma-restrict-groupoid-on-field},
l'application $p_2$ est surjective et ouverte. Ainsi, d'après
Topologie, lemme \ref{topology-lemma-open-morphism-quotient-topology},
nous voyons que $r_2'$ appartient à l'adhérence de $a'(R'_1)$.
Cela signifie que nous pouvons supposer que $r_2 \in R_2$ vérifie
que les applications $k \to \kappa(r_2)$ induites
par $s_2$ et $t_2$ sont des isomorphismes.

\medskip\noindent
Dans ce cas, nous pouvons utiliser le
lemme \ref{lemma-groupoid-on-field-move-point}.
Ce lemme implique que $c(r_2, a(R_1))$ est un voisinage ouvert de $r_2$.
Ainsi, $a(R_1) \cap c(r_2, a(R_1)) \not = \emptyset$, puisque nous avons supposé
que $r_2$ appartenait à l'adhérence de $a(R_1)$.
En utilisant les morphismes d'inversion de $R_2$ et de $R_1$, nous voyons que cela signifie que
$c_2(a(R_1), a(R_1))$ contient $r_2$.
Comme $c_2(a(R_1), a(R_1)) \subset a(c_1(R_1, R_1)) = a(R_1)$,
nous concluons que $r_2 \in a(R_1)$, comme souhaité.
\end{proof}

\begin{lemma}
\label{lemma-map-groupoids-on-field-image}
Notations et hypothèses comme dans la
situation \ref{situation-morphism-groupoids-on-field}.
Soit $Z \subset R_2$ le sous-schéma fermé réduit (voir
Schémas, définition \ref{schemes-definition-reduced-induced-scheme})
dont l'espace topologique sous-jacent est l'adhérence de l'image de
$a : R_1 \to R_2$. Alors
$c_2(Z \times_{s_2, U, t_2} Z) \subset Z$
ensemblistement.
\end{lemma}

\begin{proof}
Considérons le diagramme commutatif
$$
\xymatrix{
R_1 \times_{s_1, U, t_1} R_1 \ar[r] \ar[d] & R_1 \ar[d] \\
R_2 \times_{s_2, U, t_2} R_2 \ar[r] & R_2
}
$$
D'après
Variétés, lemme \ref{varieties-lemma-closure-image-product-map},
l'adhérence de l'image de la flèche verticale de gauche est (ensemblistement)
$Z \times_{s_2, U, t_2} Z$.
Le résultat en découle.
\end{proof}

\begin{lemma}
\label{lemma-map-groupoids-on-perfect-field-image}
Notations et hypothèses comme dans la
situation \ref{situation-morphism-groupoids-on-field}.
Supposons que $k$ soit parfait.
Soit $Z \subset R_2$ le sous-schéma fermé réduit (voir
Schémas, définition \ref{schemes-definition-reduced-induced-scheme})
dont l'espace topologique sous-jacent est l'adhérence de l'image de
$a : R_1 \to R_2$. Alors
$$
(U, Z, s_2|_Z, t_2|_Z, c_2|_Z)
$$
est un schéma en groupoïdes sur $S$.
\end{lemma}

\begin{proof}
Expliquons d'abord pourquoi l'énoncé a un sens. Puisque $U$ est le spectre
d'un corps parfait $k$, le schéma $Z$ est géométriquement réduit
sur $k$ (pour l'une ou l'autre projection), voir
Variétés, lemme \ref{varieties-lemma-perfect-reduced}.
Ainsi, le schéma $Z \times_{s_2, U, t_2} Z \subset Z$
est réduit, voir
Variétés, lemme \ref{varieties-lemma-geometrically-reduced-any-base-change}.
Par conséquent, d'après le
lemme \ref{lemma-map-groupoids-on-field-image},
nous voyons que $c$ induit un morphisme
$Z \times_{s_2, U, t_2} Z \to Z$.
Enfin, il est clair que $e_2$ se factorise par $Z$
et que le morphisme $i_2 : R_2 \to R_2$ préserve $Z$. Puisque les morphismes
du septuple
$(U, R_2, s_2, t_2, c_2, e_2, i_2)$
satisfont aux axiomes d'un groupoïde, il s'ensuit qu'après restriction
à $Z$, ils satisfont encore à ces axiomes.
\end{proof}

\begin{lemma}
\label{lemma-locally-closed-image-is-closed}
Notations et hypothèses comme dans la
situation \ref{situation-morphism-groupoids-on-field}.
Si l'image $a(R_1)$ est un sous-ensemble localement fermé de $R_2$,
alors c'est un sous-ensemble fermé.
\end{lemma}

\begin{proof}
Soit $k \subset k'$ une clôture parfaite du corps $k$.
Soit $R_i'$ la restriction de $R_i$ par le morphisme
$U' = \Spec(k') \to \Spec(k)$. Remarquons que les
morphismes $R_i' \to R_i$ sont des homéomorphismes universels, car ce sont des
composés de changements de base de l'homéomorphisme universel
$U' \to U$ (voir le diagramme dans l'énoncé du
lemme \ref{lemma-restrict-groupoid-on-field}).
Il suffit donc de démontrer que l'image $a'(R_1')$ est fermée
dans $R_2'$. Autrement dit, nous pouvons supposer que $k$ est parfait.

\medskip\noindent
Si $k$ est parfait, l'adhérence de l'image est
un schéma en groupoïdes $Z \subset R_2$, d'après le
lemme \ref{lemma-map-groupoids-on-perfect-field-image}.
En appliquant le même lemme à
$\text{id}_{R_1} : R_1 \to R_1$,
nous voyons que $(R_2)_{red}$ est un schéma en groupoïdes.
Nous pouvons donc appliquer le
lemme \ref{lemma-open-image-is-closed}
au morphisme
$a|_{(R_2)_{red}} : (R_2)_{red} \to Z$
pour conclure que $Z$ est égal à l'image de $a$.
\end{proof}

\begin{lemma}
\label{lemma-quasi-compact-map-groupoids-on-field-image}
Notations et hypothèses comme dans la
situation \ref{situation-morphism-groupoids-on-field}.
Supposons que $a : R_1 \to R_2$ soit un morphisme quasi-compact.
Soit $Z \subset R_2$ l'image schématique (voir
Morphismes, définition \ref{morphisms-definition-scheme-theoretic-image})
de $a : R_1 \to R_2$. Alors
$$
(U, Z, s_2|_Z, t_2|_Z, c_2|_Z)
$$
est un schéma en groupoïdes sur $S$.
\end{lemma}

\begin{proof}
La principale difficulté consiste à montrer que $c_2|_{Z \times_{s_2, U, t_2} Z}$
a son image dans $Z$. Considérons le diagramme commutatif
$$
\xymatrix{
R_1 \times_{s_1, U, t_1} R_1 \ar[r] \ar[d]^{a \times a} & R_1 \ar[d] \\
R_2 \times_{s_2, U, t_2} R_2 \ar[r] & R_2
}
$$
D'après le
lemme \ref{varieties-lemma-scheme-theoretic-image-product-map} du chapitre Variétés,
nous voyons que l'image schématique de $a \times a$ est
$Z \times_{s_2, U, t_2} Z$. Par commutativité du diagramme, nous
concluons que $Z \times_{s_2, U, t_2} Z$ a son image dans $Z$ par la flèche
horizontale inférieure. Comme dans la démonstration du
lemme \ref{lemma-map-groupoids-on-perfect-field-image},
on a aussi $i_2(Z) \subset Z$ et
$e_2$ se factorise par $Z$. Nous concluons donc comme dans la
démonstration de ce lemme.
\end{proof}

\begin{lemma}
\label{lemma-groupoid-on-field-image}
Soit $S$ un schéma. Soit $(U, R, s, t, c)$ un schéma en groupoïdes
sur $S$. Supposons que $U$ soit le spectre d'un corps.
Soit $Z \subset U \times_S U$ le sous-schéma fermé réduit (voir
Schémas, définition \ref{schemes-definition-reduced-induced-scheme})
dont l'espace topologique sous-jacent est l'adhérence de l'image de
$j = (t, s) : R \to U \times_S U$. Alors
$\text{pr}_{02}(Z \times_{\text{pr}_1, U, \text{pr}_0} Z) \subset Z$
ensemblistement.
\end{lemma}

\begin{proof}
Puisque $(U, U \times_S U, \text{pr}_1, \text{pr}_0, \text{pr}_{02})$
est un schéma en groupoïdes sur $S$, c'est un cas particulier du
lemme \ref{lemma-map-groupoids-on-field-image}.
Mais nous pouvons aussi le démontrer directement comme suit.

\medskip\noindent
Écrivons $U = \Spec(k)$. Notons
$R_s$ (resp.\ $Z_s$, resp.\ $U^2_s$) le schéma
$R$ (resp.\ $Z$, resp.\ $U \times_S U$) considéré comme un schéma sur $k$ au moyen de
$s$ (resp.\ $\text{pr}_1|_Z$, resp.\ $\text{pr}_1$).
De même, notons
${}_tR$ (resp.\ ${}_tZ$, resp.\ ${}_tU^2$) le schéma
$R$ (resp.\ $Z$, resp.\ $U \times_S U$) considéré comme un schéma sur $k$ au moyen de
$t$ (resp.\ $\text{pr}_0|_Z$, resp.\ $\text{pr}_0$).
Le morphisme $j$ induit des morphismes de schémas
$j_s : R_s \to U^2_s$ et ${}_tj : {}_tR \to {}_tU^2$ sur $k$.
Considérons le diagramme commutatif
$$
\xymatrix{
R_s \times_k {}_tR \ar[r]^c \ar[d]_{j_s \times {}_tj} &  R \ar[d]^j \\
U^2_s \times_k {}_tU^2 \ar[r] & U \times_S U
}
$$
D'après le
lemme \ref{varieties-lemma-closure-image-product-map} du chapitre Variétés,
nous voyons que l'adhérence de l'image de $j_s \times {}_tj$ est
$Z_s \times_k {}_tZ$. Par commutativité du diagramme, nous
concluons que $Z_s \times_k {}_tZ$ a son image dans $Z$ par la flèche
horizontale inférieure.
\end{proof}

\begin{lemma}
\label{lemma-groupoid-on-perfect-field-image}
Soit $S$ un schéma. Soit $(U, R, s, t, c)$ un schéma en groupoïdes
sur $S$. Supposons que $U$ soit le spectre d'un corps parfait.
Soit $Z \subset U \times_S U$ le sous-schéma fermé réduit (voir
Schémas, définition \ref{schemes-definition-reduced-induced-scheme})
dont l'espace topologique sous-jacent est l'adhérence de l'image de
$j = (t, s) : R \to U \times_S U$.
Alors
$$
(U, Z, \text{pr}_0|_Z, \text{pr}_1|_Z,
\text{pr}_{02}|_{Z \times_{\text{pr}_1, U, \text{pr}_0} Z})
$$
est un schéma en groupoïdes sur $S$.
\end{lemma}

\begin{proof}
Puisque $(U, U \times_S U, \text{pr}_1, \text{pr}_0, \text{pr}_{02})$
est un schéma en groupoïdes sur $S$, c'est un cas particulier du
lemme \ref{lemma-map-groupoids-on-perfect-field-image}.
Mais nous pouvons aussi le démontrer directement comme suit.

\medskip\noindent
Expliquons d'abord pourquoi l'énoncé a un sens. Puisque $U$ est le spectre
d'un corps parfait $k$, le schéma $Z$ est géométriquement réduit
sur $k$ (au moyen de l'une ou l'autre projection), voir
Variétés, lemme \ref{varieties-lemma-perfect-reduced}.
Le schéma $Z \times_{\text{pr}_1, U, \text{pr}_0} Z \subset Z$ est donc
réduit, voir
Variétés, lemme \ref{varieties-lemma-geometrically-reduced-any-base-change}.
Ainsi, d'après le
lemme \ref{lemma-groupoid-on-field-image},
nous voyons que $\text{pr}_{02}$ induit un morphisme
$Z \times_{\text{pr}_1, U, \text{pr}_0} Z \to Z$.
Enfin, il est clair que $\Delta_{U/S}$ se factorise par $Z$
et que l'application
$\sigma : U \times_S U \to U \times_S U$, $(x, y) \mapsto (y, x)$
préserve $Z$. Puisque
$(U, U \times_S U, \text{pr}_0, \text{pr}_1, \text{pr}_{02},
\Delta_{U/S}, \sigma)$
satisfait aux axiomes d'un groupoïde, il s'ensuit qu'après restriction
à $Z$, ces morphismes satisfont aux axiomes.
\end{proof}

\begin{lemma}
\label{lemma-quasi-compact-groupoid-on-field-image}
Soit $S$ un schéma. Soit $(U, R, s, t, c)$ un schéma en groupoïdes
sur $S$. Supposons que $U$ soit le spectre d'un corps et
supposons que $R$ soit quasi-compact (de manière équivalente, $s, t$ sont quasi-compacts).
Soit $Z \subset U \times_S U$ l'image schématique (voir
Morphismes, définition \ref{morphisms-definition-scheme-theoretic-image})
de $j = (t, s) : R \to U \times_S U$.
Alors
$$
(U, Z, \text{pr}_0|_Z, \text{pr}_1|_Z,
\text{pr}_{02}|_{Z \times_{\text{pr}_1, U, \text{pr}_0} Z})
$$
est un schéma en groupoïdes sur $S$.
\end{lemma}

\begin{proof}
Puisque $(U, U \times_S U, \text{pr}_1, \text{pr}_0, \text{pr}_{02})$
est un schéma en groupoïdes sur $S$, c'est un cas particulier du
lemme \ref{lemma-quasi-compact-map-groupoids-on-field-image}.
Mais nous pouvons aussi le démontrer directement comme suit.

\medskip\noindent
La difficulté principale consiste à montrer que
$\text{pr}_{02}|_{Z \times_{\text{pr}_1, U, \text{pr}_0} Z}$
a son image dans $Z$.
Écrivons $U = \Spec(k)$. Notons
$R_s$ (resp.\ $Z_s$, resp.\ $U^2_s$) le schéma
$R$ (resp.\ $Z$, resp.\ $U \times_S U$) considéré comme un schéma sur $k$ au moyen de
$s$ (resp.\ $\text{pr}_1|_Z$, resp.\ $\text{pr}_1$).
De même, notons
${}_tR$ (resp.\ ${}_tZ$, resp.\ ${}_tU^2$) le schéma
$R$ (resp.\ $Z$, resp.\ $U \times_S U$) considéré comme un schéma sur $k$ au moyen de
$t$ (resp.\ $\text{pr}_0|_Z$, resp.\ $\text{pr}_0$).
Le morphisme $j$ induit des morphismes de schémas
$j_s : R_s \to U^2_s$ et ${}_tj : {}_tR \to {}_tU^2$ sur $k$.
Considérons le diagramme commutatif
$$
\xymatrix{
R_s \times_k {}_tR \ar[r]^c \ar[d]_{j_s \times {}_tj} &  R \ar[d]^j \\
U^2_s \times_k {}_tU^2 \ar[r] & U \times_S U
}
$$
D'après le
lemme \ref{varieties-lemma-scheme-theoretic-image-product-map} du chapitre Variétés,
nous voyons que l'image schématique de $j_s \times {}_tj$ est
$Z_s \times_k {}_tZ$. Par commutativité du diagramme, nous
concluons que $Z_s \times_k {}_tZ$ a son image dans $Z$ par la flèche
horizontale inférieure. Comme dans la démonstration du
lemme \ref{lemma-groupoid-on-perfect-field-image},
on a aussi $\sigma(Z) \subset Z$ et
$\Delta_{U/S}$ se factorise par $Z$. Nous concluons donc comme dans la
démonstration de ce lemme.
\end{proof}







\section{Tranches de schémas en groupoïdes}
\label{section-slicing}

\noindent
Le lemme suivant montre que nous pouvons prendre une tranche d'un schéma en groupoïdes de Cohen-Macaulay
afin de réduire la dimension des fibres, pourvu que la
dimension du stabilisateur soit petite. C'est une étape essentielle du
processus d'amélioration d'une présentation donnée d'un champ quotient.

\begin{situation}
\label{situation-slice}
Soit $S$ un schéma.
Soit $(U, R, s, t, c)$ un schéma en groupoïdes sur $S$.
Soit $g : U' \to U$ un morphisme de schémas.
Soit $u \in U$ un point, et soit $u' \in U'$ un point tel que
$g(u') = u$. Ces données étant fixées, notons $(U', R', s', t', c')$
la restriction de $(U, R, s, t, c)$ par le morphisme $g$.
Notons $G \to U$ le schéma en groupes stabilisateur de $R$, qui
est un sous-schéma localement fermé de $R$.
Notons $h$ la composée
$$
h = s \circ \text{pr}_1 : U' \times_{g, U, t} R \longrightarrow U.
$$
Notons $F_u = s^{-1}(u)$ (fibre schématique), et $G_u$ la
fibre schématique de $G$ au-dessus de $u$.
De même, pour $R'$, notons $F'_{u'} = (s')^{-1}(u')$.
Comme $g(u') = u$, on a
$$
F'_{u'} = h^{-1}(u) \times_{\Spec(\kappa(u))} \Spec(\kappa(u')).
$$
Le point $e(u) \in R$ peut être considéré comme un point de $G_u$ et de $F_u$, et
$e'(u')$ est un point de $R'$ (resp.\ $G'_{u'}$, resp.\ $F'_{u'}$) qui s'envoie
sur $e(u)$ dans $R$ (resp.\ $G_u$, resp.\ $F_u$).
\end{situation}

\begin{lemma}
\label{lemma-slice}
Soit $S$ un schéma.
Soit $(U, R, s, t, c, e, i)$ un schéma en groupoïdes sur $S$.
Soit $G \to U$ le schéma en groupes stabilisateur.
Supposons que $s$ et $t$ soient de Cohen-Macaulay et localement de présentation finie.
Soit $u \in U$ un point de type fini du schéma $U$, voir
Morphismes, définition \ref{morphisms-definition-finite-type-point}.
Avec les notations de la
situation \ref{situation-slice},
posons
$$
d_1 = \dim(G_u), \quad
d_2 = \dim_{e(u)}(F_u).
$$
Si $d_2 > d_1$, alors il existe un schéma affine $U'$
et un morphisme $g : U' \to U$ tels que (avec les notations de la
situation \ref{situation-slice})
\begin{enumerate}
\item $g$ est une immersion,
\item $u \in U'$,
\item $g$ est localement de présentation finie,
\item le morphisme $h : U' \times_{g, U, t} R \longrightarrow U$
est de Cohen-Macaulay au point $(u, e(u))$, et
\item on a $\dim_{e'(u)}(F'_u) = d_2 - 1$.
\end{enumerate}
\end{lemma}

\begin{proof}
Soit $\Spec(A) \subset U$ un voisinage affine de $u$
tel que $u$ corresponde à un point fermé de $U$, voir
Morphismes, lemme \ref{morphisms-lemma-identify-finite-type-points}.
Soit $\Spec(B) \subset R$ un voisinage affine de $e(u)$
dont l'image par $j$ est contenue dans l'ouvert
$\Spec(A) \times_S \Spec(A) \subset U \times_S U$.
Soit $\mathfrak m \subset A$ l'idéal maximal correspondant à $u$.
Soit $\mathfrak q \subset B$ l'idéal premier correspondant à $e(u)$.
Diagrammes :
$$
\vcenter{
\xymatrix{
B & A \ar[l]^s \\
A \ar[u]^t
}
}
\quad\text{et}\quad
\vcenter{
\xymatrix{
B_{\mathfrak q} & A_{\mathfrak m} \ar[l]^s \\
A_{\mathfrak m} \ar[u]^t
}
}
$$
Notons que les deux applications induites
$s, t : \kappa(\mathfrak m) \to \kappa(\mathfrak q)$
sont égales et sont des isomorphismes puisque $s \circ e = t \circ e = \text{id}_U$.
En particulier, nous voyons que $\mathfrak q$
est aussi un idéal maximal. Les homomorphismes d'anneaux $s, t : A \to B$ sont
de présentation finie et plats. Par hypothèse, l'anneau
$$
\mathcal{O}_{F_u, e(u)} = B_{\mathfrak q}/s(\mathfrak m)B_{\mathfrak q}
$$
est de Cohen-Macaulay de dimension $d_2$. L'égalité des dimensions résulte de
Morphismes, lemme \ref{morphisms-lemma-dimension-fibre-at-a-point}.

\medskip\noindent
Soit $R''$ la restriction de $R$ à $u = \Spec(\kappa(u))$
par le morphisme $\Spec(\kappa(u)) \to U$.
Comme $u \to U$ est localement de type fini,
nous voyons que $(\Spec(\kappa(u)), R'', s'', t'', c'')$
est un schéma en groupoïdes avec $s'', t''$ localement de type fini, voir
le lemme \ref{lemma-restrict-preserves-type}.
D'après le
lemme \ref{lemma-groupoid-on-field-dimension-equal-stabilizer},
cela implique que $\dim(G'') = \dim(R'')$. On a aussi
$\dim(R'') = \dim_{e''}(R'') = \dim(\mathcal{O}_{R'', e''})$, voir
le lemme \ref{lemma-groupoid-on-field-locally-finite-type-dimension}.
D'après
Groupoïdes, lemme \ref{groupoids-lemma-restrict-stabilizer},
on a $G'' = G_u$. Nous concluons donc que
$\dim(\mathcal{O}_{R'', e''}) = d_1$.

\medskip\noindent
En tant que schéma, $R''$ est
$$
R'' =
R \times_{(U \times_S U)}
\Big(
\Spec(\kappa(\mathfrak m)) \times_S \Spec(\kappa(\mathfrak m))
\Big)
$$
Par conséquent, un voisinage ouvert affine de $e''$ est le spectre de l'anneau
$$
B \otimes_{(A \otimes A)} (\kappa(\mathfrak m) \otimes \kappa(\mathfrak m))
=
B/s(\mathfrak m)B + t(\mathfrak m)B
$$
Nous en concluons que
$$
\mathcal{O}_{R'', e''} =
B_{\mathfrak q}/s(\mathfrak m)B_{\mathfrak q} + t(\mathfrak m)B_{\mathfrak q}
$$
et nous savons donc maintenant que cet anneau est de dimension $d_1$.

\medskip\noindent
Nous affirmons que cela implique que l'on peut trouver
un élément $f \in \mathfrak m$ tel que
$$
\dim(B_{\mathfrak q}/(s(\mathfrak m)B_{\mathfrak q} + fB_{\mathfrak q}) < d_2
$$
En effet, supposons que $\mathfrak n_j \supset s(\mathfrak m)B_{\mathfrak q}$,
$j = 1, \ldots, m$, correspondent aux idéaux premiers minimaux de l'anneau local
$B_{\mathfrak q}/s(\mathfrak m)B_{\mathfrak q}$.
Ils sont en nombre fini puisque cet anneau est noethérien (car il est essentiellement
de type fini sur un corps -- mais aussi parce qu'un anneau de Cohen-Macaulay est
noethérien). Par la propriété de Cohen-Macaulay, on a
$\dim(B_{\mathfrak q}/\mathfrak n_j) = d_2$, par exemple d'après
Algèbre, lemme \ref{algebra-lemma-CM-dim-formula}.
Notons que
$\dim(B_{\mathfrak q}/(\mathfrak n_j + t(\mathfrak m)B_{\mathfrak q}))
\leq d_1$
car c'est un quotient de l'anneau
$\mathcal{O}_{R'', e''} =
B_{\mathfrak q}/s(\mathfrak m)B_{\mathfrak q} + t(\mathfrak m)B_{\mathfrak q}$
qui est de dimension $d_1$. Comme $d_1 < d_2$, cela implique que
$\mathfrak m \not \subset t^{-1}(\mathfrak n_i)$.
Par le lemme d'évitement des idéaux premiers, voir
Algèbre, lemme \ref{algebra-lemma-silly},
nous pouvons trouver $f \in \mathfrak m$ tel que $t(f) \not \in \mathfrak n_j$ pour
$j = 1, \ldots, m$. Pour ce choix de $f$, on a
l'inégalité affichée ci-dessus, voir
Algèbre, lemme \ref{algebra-lemma-one-equation}.

\medskip\noindent
Posons $A' = A/fA$ et $U' = \Spec(A')$. Il est alors clair que
$U' \to U$ est une immersion, localement de présentation finie,
et que $u \in U'$. Ainsi, les assertions (1), (2) et (3) du lemme sont vérifiées.
Le morphisme
$$
U' \times_{g, U, t} R \longrightarrow U
$$
se factorise par $\Spec(A)$ et correspond à l'homomorphisme d'anneaux
$$
\xymatrix{
B/t(f)B \ar@{=}[r] & A/(f) \otimes_{A, t} B & A \ar[l]_-s
}
$$
Nous voyons maintenant que $t(f)$ n'est pas un diviseur de zéro dans
$B_{\mathfrak q}/s(\mathfrak m)B_{\mathfrak q}$, car c'est un anneau de
Cohen-Macaulay de dimension positive et $f$ n'appartient à
aucun idéal premier minimal, voir par exemple
Algèbre, lemme \ref{algebra-lemma-reformulate-CM}.
Par conséquent, d'après
Algèbre, lemme \ref{algebra-lemma-grothendieck-general},
nous concluons que $s : A_{\mathfrak m} \to B_{\mathfrak q}/t(f)B_{\mathfrak q}$
est plat et a pour anneau de la fibre
$B_{\mathfrak q}/(s(\mathfrak m)B_{\mathfrak q} + t(f)B_{\mathfrak q})$,
qui est de Cohen-Macaulay, encore d'après
Algèbre, lemme \ref{algebra-lemma-reformulate-CM}.
Cela implique l'assertion (4) du lemme.
Pour voir l'assertion (5), notons que, d'après le diagramme (\ref{equation-restriction}),
la fibre $F'_u$ est égale à la fibre de $h$ au-dessus de $u$.
Par conséquent,
$\dim_{e'(u)}(F'_u) =
\dim(B_{\mathfrak q}/(s(\mathfrak m)B_{\mathfrak q} + t(f)B_{\mathfrak q}))$
d'après
Morphismes, lemme \ref{morphisms-lemma-dimension-fibre-at-a-point},
et la dimension de cet anneau est $d_2 - 1$ d'après
Algèbre, lemme \ref{algebra-lemma-reformulate-CM},
une fois encore. Cela démontre la dernière assertion du lemme et achève la démonstration.
\end{proof}

\noindent
Maintenant que nous savons prendre une tranche, nous pouvons combiner ce résultat avec ce qui précède
pour obtenir le résultat ``optimal'' suivant. Il est optimal au sens où,
puisque $G_u$ est un sous-schéma localement fermé de $F_u$, on a toujours
l'inégalité $\dim(G_u) = \dim_{e(u)}(G_u) \leq \dim_{e(u)}(F_u)$ ; il
n'est donc pas possible de pousser plus loin la prise d'une tranche que dans le lemme.

\begin{lemma}
\label{lemma-max-slice}
Soit $S$ un schéma.
Soit $(U, R, s, t, c, e, i)$ un schéma en groupoïdes sur $S$.
Soit $G \to U$ le schéma en groupes stabilisateur.
Supposons que $s$ et $t$ soient de Cohen-Macaulay et localement de présentation finie.
Soit $u \in U$ un point de type fini du schéma $U$, voir
Morphismes, définition \ref{morphisms-definition-finite-type-point}.
Avec les notations de la
situation \ref{situation-slice},
il existe un schéma affine $U'$ et un morphisme $g : U' \to U$ tels que
\begin{enumerate}
\item $g$ est une immersion,
\item $u \in U'$,
\item $g$ est localement de présentation finie,
\item le morphisme $h : U' \times_{g, U, t} R \longrightarrow U$
est de Cohen-Macaulay et localement de présentation finie,
\item les morphismes $s', t' : R' \to U'$ sont de Cohen-Macaulay et
localement de présentation finie, et
\item $\dim_{e(u)}(F'_u) = \dim(G'_u)$.
\end{enumerate}
\end{lemma}

\begin{proof}
Comme $s$ est localement de présentation finie, le schéma $F_u$ est
localement de type fini sur $\kappa(u)$. Par conséquent,
$\dim_{e(u)}(F_u) < \infty$ et nous pouvons raisonner par récurrence sur
$\dim_{e(u)}(F_u)$.

\medskip\noindent
Si $\dim_{e(u)}(F_u) = \dim(G_u)$, il n'y a rien à démontrer.
Supposons $\dim_{e(u)}(F_u) > \dim(G_u)$. Cela signifie que le
lemme \ref{lemma-slice}
s'applique, et nous trouvons un morphisme $g : U' \to U$ qui possède
les propriétés (1), (2), (3), tandis qu'au lieu de (6), nous avons
$\dim_{e(u)}(F'_u) < \dim_{e(u)}(F_u)$,
et qu'au lieu de (4) et (5), nous savons que la composée
$$
h = s \circ \text{pr}_1 : U' \times_{g, U, t} R \longrightarrow U
$$
est de Cohen-Macaulay au point $(u, e(u))$. Nous appliquons la
remarque \ref{remark-local-source-apply}
et obtenons un sous-schéma ouvert $U'' \subset U'$ tel que
$U'' \times_{g, U, t} R \subset U' \times_{g, U, t} R$
soit le plus grand sous-schéma ouvert sur lequel $h$ est de Cohen-Macaulay.
Puisque $(u, e(u)) \in U'' \times_{g, U, t} R$, nous voyons que $u \in U''$.
Nous pouvons donc remplacer $U'$ par $U''$ et supposer qu'en fait $h$ est
de Cohen-Macaulay partout ! D'après le
lemme \ref{lemma-restrict-property},
nous concluons que $s', t'$ sont localement de
présentation finie et de Cohen-Macaulay (utiliser
Morphismes, lemme \ref{morphisms-lemma-base-change-finite-presentation},
et
Compléments sur les morphismes, lemme \ref{more-morphisms-lemma-base-change-CM}).

\medskip\noindent
Par construction, $\dim_{e'(u)}(F'_u) < \dim_{e(u)}(F_u)$,
de sorte que nous pouvons appliquer l'hypothèse de récurrence à $(U', R', s', t', c')$
et au point $u \in U'$. Notons que $u$ est également un point de type fini
de $U'$ (on peut par exemple le voir en utilisant la caractérisation des
points de type fini donnée dans
Morphismes, lemme \ref{morphisms-lemma-identify-finite-type-points}).
Soient $g' : U'' \to U'$ et $(U'', R'', s'', t'', c'')$ la solution
du problème correspondant partant de $(U', R', s', t', c')$
et du point $u \in U'$. Nous affirmons que la composée
$$
g'' = g \circ g' : U'' \longrightarrow U
$$
est une solution du problème initial. Les propriétés (1), (2), (3), (5)
et (6) sont immédiates. Pour voir (4), notons que le morphisme
$$
h'' = s \circ \text{pr}_1 : U'' \times_{g'', U, t} R \longrightarrow U
$$
est localement de présentation finie et de Cohen-Macaulay par application du
lemme \ref{lemma-double-restrict-property}
(utiliser
Compléments sur les morphismes, lemme \ref{more-morphisms-lemma-CM-local-source-and-target},
pour voir que la propriété d'être de Cohen-Macaulay est locale pour la topologie fppf sur la base).
\end{proof}

\noindent
Lorsque les fibres du schéma en groupes stabilisateur sont de dimension 0,
cela donne le lemme de tranche suivant.

\begin{lemma}
\label{lemma-max-slice-quasi-finite}
Soit $S$ un schéma.
Soit $(U, R, s, t, c, e, i)$ un schéma en groupoïdes sur $S$.
Soit $G \to U$ le schéma en groupes stabilisateur.
Supposons que $s$ et $t$ soient de Cohen-Macaulay et localement de présentation finie.
Soit $u \in U$ un point de type fini du schéma $U$, voir
Morphismes, définition \ref{morphisms-definition-finite-type-point}.
Supposons que $G \to U$ soit localement quasi-fini.
Avec les notations de la
situation \ref{situation-slice},
il existe un schéma affine $U'$ et un morphisme $g : U' \to U$ tels que
\begin{enumerate}
\item $g$ est une immersion,
\item $u \in U'$,
\item $g$ est localement de présentation finie,
\item le morphisme $h : U' \times_{g, U, t} R \longrightarrow U$
est plat, localement de présentation finie et localement quasi-fini, et
\item les morphismes $s', t' : R' \to U'$ sont plats,
localement de présentation finie et localement quasi-finis.
\end{enumerate}
\end{lemma}

\begin{proof}
Prenons $g : U' \to U$ comme dans le
lemme \ref{lemma-max-slice}.
Puisque $h^{-1}(u) = F'_u$, nous voyons que $h$ est de dimension relative
$\leq 0$ en $(u, e(u))$. Par conséquent, d'après la
remarque \ref{remark-local-source-apply},
nous obtenons un sous-schéma ouvert $U'' \subset U'$ tel que
$u \in U''$ et que $U'' \times_{g, U, t} R$ soit le plus grand sous-schéma ouvert
de $U' \times_{g, U, t} R$ sur lequel $h$ est de dimension relative $\leq 0$.
Après avoir remplacé $U'$ par $U''$, nous voyons que $h$ est de dimension relative $\leq 0$.
Cela implique que $h$ est localement quasi-fini d'après
Morphismes, lemme \ref{morphisms-lemma-locally-quasi-finite-rel-dimension-0}.
Comme il est encore localement de présentation finie et de Cohen-Macaulay, nous voyons
qu'il est plat, localement de présentation finie et localement quasi-fini,
c'est-à-dire que (4) ci-dessus est vérifiée. Cela implique que $s'$ est plat, localement de
présentation finie et localement quasi-fini en tant que changement de base de $h$, voir
le lemme \ref{lemma-restrict-property}.
\end{proof}







\section{Localisation \'etale des schémas en groupoïdes}
\label{section-etale-localize}

\noindent
Dans cette section, nous commençons à appliquer aux schémas en groupoïdes les techniques de localisation \'etale de
Compléments sur les morphismes, section \ref{more-morphisms-section-etale-localization}.
On trouvera des résultats plus avancés de ce type dans
Compléments sur les groupoïdes dans les espaces,
section \ref{spaces-more-groupoids-section-etale-localize}.
Le lemme \ref{lemma-quasi-finite-over-base-j-proper}
sera utilisé pour démontrer des résultats sur les espaces algébriques
séparés et quasi-finis sur un schéma, à savoir
Morphismes d'espaces, proposition
\ref{spaces-morphisms-proposition-locally-quasi-finite-separated-over-scheme}
et son corollaire
Morphismes d'espaces, lemme
\ref{spaces-morphisms-lemma-locally-quasi-finite-separated-representable}.

\begin{lemma}
\label{lemma-quasi-finite-over-base}
Soit $S$ un schéma.
Soit $(U, R, s, t, c)$ un schéma en groupoïdes sur $S$.
Soit $p \in S$ un point, et soit $u \in U$ un point situé au-dessus de $p$.
Supposons que
\begin{enumerate}
\item $U \to S$ soit localement de type fini,
\item $U \to S$ soit quasi-fini en $u$,
\item $U \to S$ soit séparé,
\item $R \to S$ soit séparé,
\item $s$, $t$ soient plats et localement de présentation finie, et
\item $s^{-1}(\{u\})$ soit fini.
\end{enumerate}
Alors il existe un voisinage \'etale $(S', p') \to (S, p)$ tel que
$\kappa(p) = \kappa(p')$ et un diagramme de changement de base
$$
\xymatrix{
R' \amalg W'
\ar@{=}[r] &
S' \times_S R
\ar[r] \ar@<2ex>[d]^{s'} \ar@<-2ex>[d]_{t'} &
R \ar@<1ex>[d]^s \ar@<-1ex>[d]_t \\
U' \amalg W
\ar@{=}[r] &
S' \times_S U
\ar[r] \ar[d] &
U \ar[d] \\
 &
S' \ar[r] &
S
}
$$
où les signes d'égalité sont des décompositions en sous-schémas ouverts et
fermés telles que
\begin{enumerate}
\item[(a)] il existe un point $u'$ de $U'$ qui s'envoie sur $u$ dans $U$,
\item[(b)] la fibre $(U')_{p'}$ est égale à $t'\big((s')^{-1}(\{u'\})\big)$
ensemblistement,
\item[(c)] la fibre $(R')_{p'}$ est égale à $(s')^{-1}\big((U')_{p'}\big)$
ensemblistement,
\item[(d)] les schémas $U'$ et $R'$ sont finis sur $S'$,
\item[(e)] on a $s'(R') \subset U'$ et $t'(R') \subset U'$,
\item[(f)] on a
$c'(R' \times_{s', U', t'} R') \subset R'$
où $c'$ est le changement de base de $c$, et
\item[(g)] les morphismes $s', t', c'$ déterminent une structure de groupoïde
en prenant le système
$(U', R', s'|_{R'}, t'|_{R'}, c'|_{R' \times_{s', U', t'} R'})$.
\end{enumerate}
\end{lemma}

\begin{proof}
Notons $f : U \to S$ le morphisme structural de $U$.
D'après l'hypothèse (6), on peut écrire $s^{-1}(\{u\}) = \{r_1, \ldots, r_n\}$.
Comme cet ensemble est fini, on voit que $s$ est quasi-fini en chacun de
ces antécédents, en nombre fini, voir
Morphismes, Lemme \ref{morphisms-lemma-finite-fibre}.
On voit donc que $f \circ s : R \to S$ est quasi-fini en chaque $r_i$
(Morphismes, Lemme \ref{morphisms-lemma-composition-quasi-finite}).
Ainsi, $r_i$ est isolé dans la fibre $R_p$, voir
Morphismes, Lemme \ref{morphisms-lemma-quasi-finite-at-point-characterize}.
Écrivons $t(\{r_1, \ldots, r_n\}) = \{u_1, \ldots, u_m\}$.
Notons qu'il peut arriver que $m < n$ et que
$u \in \{u_1, \ldots, u_m\}$.
Comme $t$ est plat et localement de présentation finie,
le morphisme de fibres $t_p : R_p \to U_p$ est plat et localement de
présentation finie (Morphismes,
Lemmes \ref{morphisms-lemma-base-change-flat} et
\ref{morphisms-lemma-base-change-finite-presentation}),
donc ouvert (Morphismes,
Lemme \ref{morphisms-lemma-fppf-open}).
Le fait que chaque $r_i$ soit isolé dans $R_p$ implique que
chaque $u_j = t(r_i)$ est isolé dans $U_p$. En utilisant de nouveau
Morphismes, Lemme \ref{morphisms-lemma-quasi-finite-at-point-characterize},
on voit que $f$ est quasi-fini en $u_1, \ldots, u_m$.

\medskip\noindent
Notons $F_u = s^{-1}(u)$ et $F_{u_j} = s^{-1}(u_j)$ les fibres schématiques.
Notons que $F_u$ est fini sur $\kappa(u)$, car il est localement de type fini
sur $\kappa(u)$ et possède un nombre fini de points (cela résulte par exemple
du résultat beaucoup plus général
Morphismes, Lemme
\ref{morphisms-lemma-locally-quasi-finite-qc-source-universally-bounded}).
D'après
le Lemme \ref{lemma-two-fibres},
on voit que $F_u$ et $F_{u_j}$ deviennent isomorphes après extension à un corps
qui est une extension commune de $\kappa(u)$ et $\kappa(u_j)$. On voit donc
que $F_{u_j}$ est fini sur $\kappa(u_j)$. En particulier, on voit que
$s^{-1}(\{u_j\})$ est un ensemble fini pour chaque $j = 1, \ldots, m$.
Ainsi, les hypothèses (2) et (6) sont aussi vérifiées pour chaque $u_j$
(on a vu plus haut que $U \to S$ est quasi-fini en $u_j$).
L'argument du premier paragraphe s'applique donc à chaque $u_j$,
et l'on voit que $R \to U$ est quasi-fini en chacun des points de
$$
\{r_1, \ldots, r_N\} = s^{-1}(\{u_1, \ldots, u_m\})
$$
Notons que $t(\{r_1, \ldots, r_N\}) = \{u_1, \ldots, u_m\}$ et
$t^{-1}(\{u_1, \ldots, u_m\}) = \{r_1, \ldots, r_N\}$
car $R$ est un groupoïde\footnote{Explication dans le langage des groupoïdes :
L'ensemble initial $\{r_1, \ldots, r_n\}$ était l'ensemble des flèches de
source $u$. L'ensemble $\{u_1, \ldots, u_m\}$ était l'ensemble des objets
isomorphes à $u$. Et $\{r_1, \ldots, r_N\}$ est l'ensemble de toutes les flèches
entre tous les objets équivalents à $u$.}. De plus, on a
$\text{pr}_0(c^{-1}(\{r_1, \ldots, r_N\})) = \{r_1, \ldots, r_N\}$
et
$\text{pr}_1(c^{-1}(\{r_1, \ldots, r_N\})) = \{r_1, \ldots, r_N\}$.
De même, on obtient $e(\{u_1, \ldots, u_m\}) \subset \{r_1, \ldots, r_N\}$
et $i(\{r_1, \ldots, r_N\}) = \{r_1, \ldots, r_N\}$.

\medskip\noindent
On peut appliquer
Compléments sur les morphismes,
Lemme \ref{more-morphisms-lemma-etale-splits-off-quasi-finite-part-technical}
aux couples
$(U \to S, \{u_1, \ldots, u_m\})$ et
$(R \to S, \{r_1, \ldots, r_N\})$
pour obtenir un voisinage \'etale $(S', p') \to (S, p)$
qui induit une identification $\kappa(p) = \kappa(p')$
tel que $S' \times_S U$ et $S' \times_S R$ se décomposent sous la forme
$$
S' \times_S U = U' \amalg W, \quad
S' \times_S R = R' \amalg W'
$$
où $U' \to S'$ est fini et $(U')_{p'}$ s'envoie bijectivement sur
$\{u_1, \ldots, u_m\}$, et où $R' \to S'$ est fini et
$(R')_{p'}$ s'envoie bijectivement sur $\{r_1, \ldots, r_N\}$.
De plus, aucun point de $W_{p'}$ (resp.\ de $(W')_{p'}$) n'a pour image
aucun des points $u_j$ (resp.\ $r_i$). À ce stade, (a), (b), (c) et (d)
du lemme sont satisfaites. En outre, les inclusions de (e) et (f) valent
sur les fibres au-dessus de $p'$, c'est-à-dire $s'((R')_{p'}) \subset (U')_{p'}$,
$t'((R')_{p'}) \subset (U')_{p'}$, et
$c'((R' \times_{s', U', t'} R')_{p'}) \subset (R')_{p'}$.

\medskip\noindent
Nous affirmons que l'on peut remplacer $S'$ par un voisinage ouvert de Zariski
de $p'$ de sorte que les inclusions de (e) et (f) soient vérifiées.
Considérons par exemple l'ensemble $E = (s'|_{R'})^{-1}(W)$.
Il est ouvert et fermé dans $R'$ et ne contient aucun point
de $R'$ situé au-dessus de $p'$. Comme $R' \to S'$ est fermé,
après avoir remplacé $S'$ par $S' \setminus (R' \to S')(E)$, on arrive à une
situation où $E$ est vide. Autrement dit, $s'$ envoie $R'$ dans $U'$.
Notons que cette propriété est conservée lorsque l'on rétrécit encore $S'$.
De même, on peut faire en sorte que $t'$ envoie $R'$ dans $U'$.
À ce stade, (e) est vérifiée. De la même manière, considérons l'ensemble
$E = (c'|_{R' \times_{s', U', t'} R'})^{-1}(W')$.
Il est ouvert et fermé dans le schéma $R' \times_{s', U', t'} R'$,
qui est fini sur $S'$, et ne contient aucun point situé
au-dessus de $p'$. Après avoir remplacé $S'$ par
$S' \setminus (R' \times_{s', U', t'} R' \to S')(E)$,
on arrive donc à une situation où $E$ est vide. Autrement dit, on obtient
l'inclusion de (f). On peut répéter l'argument avec l'identité
$e' : S' \times_S U \to S' \times_S R$ et l'inverse
$i' : S' \times_S R \to S' \times_S R$, afin de pouvoir supposer
(après avoir encore rétréci $S'$) que $(e'|_{U'})^{-1}(W') = \emptyset$
et $(i'|_{R'})^{-1}(W') = \emptyset$.

\medskip\noindent
À ce stade, on voit que l'on peut considérer la structure
$$
(U', R', s'|_{R'}, t'|_{R'}, c'|_{R' \times_{t', U', s'} R'},
e'|_{U'}, i'|_{R'}).
$$
Les axiomes d'un schéma en groupoïdes sur $S'$ sont vérifiés
parce qu'ils le sont pour le schéma en groupoïdes
$(S' \times_S U, S' \times_S R, s', t', c', e', i')$.
\end{proof}

\begin{lemma}
\label{lemma-quasi-finite-over-base-j-proper}
Soit $S$ un schéma.
Soit $(U, R, s, t, c)$ un schéma en groupoïdes sur $S$.
Soient $p \in S$ un point et $u \in U$ un point situé au-dessus de $p$.
Supposons que les hypothèses (1) -- (6) du
Lemme \ref{lemma-quasi-finite-over-base}
soient vérifiées, ainsi que
\begin{enumerate}
\item[(7)] $j : R \to U \times_S U$ est universellement fermé\footnote{Compte tenu
des autres conditions, cela équivaut à demander que $j$ soit propre.}.
\end{enumerate}
Alors on peut choisir $(S', p') \to (S, p)$, des décompositions
$S' \times_S U = U' \amalg W$ et $S' \times_S R = R' \amalg W'$
et $u' \in U'$ tels que (a) -- (g) du
Lemme \ref{lemma-quasi-finite-over-base}
soient vérifiées, ainsi que
\begin{enumerate}
\item[(h)] $R'$ est la restriction de $S' \times_S R$ à $U'$.
\end{enumerate}
\end{lemma}

\begin{proof}
On applique le Lemme \ref{lemma-quasi-finite-over-base} au
groupoïde $(U, R, s, t, c)$ sur le schéma $S$, avec les points $p$ et $u$.
On obtient donc un voisinage \'etale
$(S', p') \to (S, p)$ et des décompositions en unions disjointes
$$
S' \times_S U = U' \amalg W, \quad
S' \times_S R = R' \amalg W'
$$
ainsi que $u' \in U'$ satisfaisant les conclusions (a), (b), (c), (d), (e), (f) et (g).
On peut rétrécir $S'$ à un voisinage plus petit de $p'$ sans
affecter les conclusions (a) -- (g). Nous allons montrer que, pour un
rétrécissement convenable, la conclusion (h) est également vérifiée.
Notons $j'$ le changement de base de $j$ à $S'$.
D'après la conclusion (e), il est clair que
$$
j'^{-1}(U' \times_{S'} U') = R' \amalg Rest
$$
pour une certaine partie $Rest$ ouverte et fermée. Comme $U' \to S'$ est fini
d'après la conclusion (d), on voit que $U' \times_{S'} U'$ est fini sur $S'$.
Comme $j$ est universellement fermé, $j'$ l'est aussi, et
donc $j'|_{Rest}$ est lui aussi universellement fermé. D'après les conclusions
(b) et (c), on voit que la fibre de
$$
(U' \times_{S'} U' \to S') \circ j'|_{Rest} :
Rest
\longrightarrow
S'
$$
au-dessus de $p'$ est vide. Ainsi, puisque $Rest \to S'$ est fermé comme composé
de morphismes fermés, après avoir remplacé $S'$ par
$S' \setminus \Im(Rest \to S')$, on peut supposer que
$Rest = \emptyset$. C'est exactement la condition disant que $R'$ est
la restriction de $S' \times_S R$ au sous-schéma ouvert
$U' \subset S' \times_S U$, voir
Groupoïdes, Lemme \ref{groupoids-lemma-restrict-groupoid-relation}
et sa démonstration.
\end{proof}






\section{Groupoïdes finis}
\label{section-finite-groupoids}

\noindent
Un schéma en groupoïdes $(U, R, s, t, c)$ est parfois dit {\it fini} si les
morphismes $s$ et $t$ sont finis. Cette terminologie peut prêter à confusion, car elle
n'implique pas que $U$, $R$ ou le faisceau quotient $U/R$ soient finis sur quoi que ce soit.

\begin{lemma}
\label{lemma-finite-stratify}
Soit $(U, R, s, t, c)$ un schéma en groupoïdes sur un schéma $S$. Supposons $s, t$
finis. Il existe une suite de sous-schémas fermés invariants par $R$
$$
U = Z_0 \supset Z_1 \supset Z_2 \supset \ldots
$$
telle que $\bigcap Z_r = \emptyset$ et que
$s^{-1}(Z_{r - 1}) \setminus s^{-1}(Z_r) \to Z_{r - 1} \setminus Z_r$
soit fini localement libre de rang $r$.
\end{lemma}

\begin{proof}
Soit $\{Z_r\}$ la stratification de $U$ donnée par les idéaux de Fitting
du module quasi-cohérent de type fini $s_*\mathcal{O}_R$. Voir
Diviseurs, Lemme \ref{divisors-lemma-locally-free-rank-r-pullback}.
Comme l'identité $e : U \to R$ est une section de $s$, on voit que
$s_*\mathcal{O}_R$ contient $\mathcal{O}_S$ comme facteur direct.
Ainsi $U = Z_{-1} = Z_0$ (détails omis).
Comme la formation des idéaux de Fitting commute au changement de base
(Compléments sur l'algèbre, Lemme \ref{more-algebra-lemma-fitting-ideal-basics}),
on trouve que $s^{-1}(Z_r)$ correspond au $r$-ième idéal de Fitting
de $\text{pr}_{1, *}\mathcal{O}_{R \times_{s, U, t} R}$, car
le carré inférieur droit du diagramme (\ref{equation-diagram}) est cartésien.
En utilisant le fait que le carré inférieur gauche est lui aussi cartésien, on conclut
que $s^{-1}(Z_r) = t^{-1}(Z_r)$ ; autrement dit, $Z_r$ est invariant par $R$.
Le morphisme
$s^{-1}(Z_{r - 1}) \setminus s^{-1}(Z_r) \to Z_{r - 1} \setminus Z_r$
est fini localement libre de rang $r$, car le module
$s_*\mathcal{O}_R$ se restreint en un module fini localement libre de rang $r$
sur $Z_{r - 1} \setminus Z_r$, d'après
Diviseurs, Lemme \ref{divisors-lemma-locally-free-rank-r-pullback}.
\end{proof}

\begin{lemma}
\label{lemma-finite-flat-over-almost-dense-subscheme}
Soit $(U, R, s, t, c)$ un schéma en groupoïdes sur un schéma $S$. Supposons $s, t$
finis. Il existe un sous-schéma ouvert $W \subset U$ et un sous-schéma fermé
$W' \subset W$ tels que
\begin{enumerate}
\item $W$ et $W'$ soient invariants par $R$,
\item $U = t(s^{-1}(\overline{W}))$ ensemblistement,
\item $W$ soit un épaississement de $W'$, et
\item les morphismes $s'$, $t'$ de la restriction $(W', R', s', t', c')$
soient finis localement libres.
\end{enumerate}
\end{lemma}

\begin{proof}
Considérons la stratification $U = Z_0 \supset Z_1 \supset Z_2 \supset \ldots$
du Lemme \ref{lemma-finite-stratify}.

\medskip\noindent
Nous allons construire des unions disjointes $W = \coprod_{r \geq 1} W_r$ et
$W' = \coprod_{r \geq 1} W'_r$, chaque $W'_r \to W_r$ étant un épaississement
entre des sous-schémas invariants par $R$ de $U$, de sorte que les morphismes
$s_r', t_r'$ des restrictions $(W_r', R_r', s_r', t_r', c_r')$
soient finis localement libres de rang $r$. Pour commencer, posons
$W_1 = W'_1 = U \setminus Z_1$. C'est un sous-schéma ouvert invariant par $R$
de $U$ ; il est vrai que $W_0$ est un épaississement de $W'_0$,
et les morphismes $s_1'$, $t_1'$ de la
restriction $(W_1', R_1', s_1', t_1', c_1')$ sont des isomorphismes, c'est-à-dire
finis localement libres de rang $1$.
De plus, tout point de $U \setminus Z_1$ appartient à $t(s^{-1}(\overline{W_1}))$.

\medskip\noindent
Supposons avoir trouvé des sous-schémas $W'_r \subset W_r \subset U$ pour $r \leq n$
tels que
\begin{enumerate}
\item $W_1, \ldots, W_n$ soient disjoints,
\item $W_r$ et $W_r'$ soient invariants par $R$,
\item $U \setminus Z_n \subset \bigcup_{r \leq n} t(s^{-1}(\overline{W_r}))$
ensemblistement,
\item $W_r$ soit un épaississement de $W'_r$,
\item les morphismes $s_r'$, $t_r'$ de la restriction
$(W_r', R_r', s_r', t_r', c_r')$ soient finis localement libres de rang $r$.
\end{enumerate}
Posons alors
$$
W_{n + 1} = Z_n \setminus
\left(
Z_{n + 1} \cup \bigcup\nolimits_{r \leq n} t(s^{-1}(\overline{W_r}))
\right)
$$
ensemblistement et
$$
W'_{n + 1} = Z_n \setminus
\left(
Z_{n + 1} \cup \bigcup\nolimits_{r \leq n} t(s^{-1}(\overline{W_r}))
\right)
$$
schématiquement. Alors $W_{n + 1}$ est un sous-schéma ouvert invariant par $R$
de $U$, car $Z_{n + 1} \setminus \overline{U \setminus Z_{n + 1}}$
est ouvert dans $U$ et $\overline{U \setminus Z_{n + 1}}$ est contenu
dans le fermé $\bigcup\nolimits_{r \leq n} t(s^{-1}(\overline{W_r}))$
que l'on retire, par la propriété (3) et le fait que $t$ est un morphisme fermé.
Il est clair que $W'_{n + 1}$ est un sous-schéma fermé
de $W_{n + 1}$ ayant le même espace topologique sous-jacent.
Enfin, les propriétés (1), (2) et (3) sont claires, et la propriété (5) résulte du
Lemme \ref{lemma-finite-stratify}.

\medskip\noindent
D'après le Lemme \ref{lemma-finite-stratify}, on a $\bigcap Z_r = \emptyset$.
Tout point de $U$ est donc contenu dans $U \setminus Z_n$
pour un certain $n$. On voit ainsi que
$U = \bigcup_{r \geq 1} t(s^{-1}(\overline{W_r}))$
ensemblistement, et l'on voit que (2) est vérifiée.
Ainsi, les sous-schémas $W' \subset W$ satisfont (1), (2), (3) et (4).
\end{proof}

\noindent
Soit $(U, R, s, t, c)$ un schéma en groupoïdes. Étant donné un point $u \in U$,
l'{\it orbite sous $R$} de $u$ est le sous-ensemble $t(s^{-1}(\{u\}))$ de $U$.

\begin{lemma}
\label{lemma-finite-flat-over-almost-dense-subscheme-addendum}
Dans le Lemme \ref{lemma-finite-flat-over-almost-dense-subscheme},
supposons en outre que $s$ et $t$ soient de présentation finie.
Alors
\begin{enumerate}
\item le morphisme $W' \to W$ est de présentation finie, et
\item si $u \in U$ est un point dont l'orbite sous $R$ est constituée de
points génériques de composantes irréductibles de $U$, alors $u \in W$.
\end{enumerate}
\end{lemma}

\begin{proof}
Dans ce cas, la stratification
$U = Z_0 \supset Z_1 \supset Z_2 \supset \ldots$ du
Lemme \ref{lemma-finite-stratify} est donnée par des immersions fermées $Z_k \to U$
de présentation finie, voir
Diviseurs, Lemme \ref{divisors-lemma-locally-free-rank-r-pullback}.
La partie (1) en résulte immédiatement, puisque $W' \to W$ est localement obtenu
par l'intersection de l'ouvert $W$ avec $Z_r$. Pour voir la partie (2),
soit $\{u_1, \ldots, u_n\}$ l'orbite de $u$.
Comme les sous-schémas fermés $Z_k$ sont invariants par $R$ et que
$\bigcap Z_k = \emptyset$, on trouve un $k$ tel que $u_i \in Z_k$
et $u_i \not \in Z_{k + 1}$ pour tout $i$.
L'image de $Z_k \to U$ et celle de $Z_{k + 1} \to U$ sont localement constructibles
(Morphismes, Théorème \ref{morphisms-theorem-chevalley}).
Comme $u_i \in U$ est un point générique d'une composante irréductible
de $U$, il existe un voisinage ouvert $U_i$ de $u_i$ qui
est contenu dans $Z_k \setminus Z_{k + 1}$ ensemblistement
(Propriétés, Lemme \ref{properties-lemma-generic-point-in-constructible}).
Dans la démonstration du Lemme \ref{lemma-finite-flat-over-almost-dense-subscheme},
nous avons construit $W$ comme une union disjointe $\coprod W_r$,
avec $W_r \subset Z_{r - 1} \setminus Z_r$, de sorte que
$U = \bigcup t(s^{-1}(\overline{W_r}))$. Comme $\{u_1, \ldots, u_n\}$
est une orbite sous $R$, on voit que $u \in t(s^{-1}(\overline{W_r}))$
implique $u_i \in \overline{W_r}$ pour un certain $i$, ce qui implique
$U_i \cap W_r \not = \emptyset$, et donc $r = k$.
On en conclut que $u$ appartient à
$$
W_{k + 1} = Z_k \setminus
\left(
Z_{k + 1} \cup \bigcup\nolimits_{r \leq k} t(s^{-1}(\overline{W_r}))
\right)
$$
comme voulu.
\end{proof}

\begin{lemma}
\label{lemma-invariant-affine-open-around-generic-point}
Soit $(U, R, s, t, c)$ un schéma en groupoïdes sur un schéma $S$. Supposons $s, t$
finis et de présentation finie, et $U$ quasi-séparé. Soient
$u_1, \ldots, u_m \in U$ des points dont les orbites sont constituées de points génériques
de composantes irréductibles de $U$. Il existe alors des sous-schémas invariants par $R$
$V' \subset V \subset U$ tels que
\begin{enumerate}
\item $u_1, \ldots, u_m \in V'$,
\item $V$ soit ouvert dans $U$,
\item $V'$ et $V$ soient affines,
\item $V' \subset V$ soit un épaississement de présentation finie,
\item les morphismes $s', t'$ de la restriction $(V', R', s', t', c')$
soient finis localement libres.
\end{enumerate}
\end{lemma}

\begin{proof}
Soient $W' \subset W \subset U$ comme dans le
Lemme \ref{lemma-finite-flat-over-almost-dense-subscheme}.
D'après le Lemme \ref{lemma-finite-flat-over-almost-dense-subscheme-addendum},
on a $u_j \in W$ pour chaque indice, et $W' \to W$ est un épaississement de présentation finie.
D'après Limites, Lemme \ref{limits-lemma-affines-glued-in-closed-affine},
il suffit de trouver un sous-schéma ouvert affine invariant par $R$
$V'$ de $W'$ contenant $u_j$ (car on peut alors prendre pour $V \subset W$
le sous-schéma ouvert correspondant, qui sera affine).
On peut donc remplacer $(U, R, s, t, c)$ par la restriction
$(W', R', s', t', c')$ à $W'$.
Autrement dit, on peut supposer que l'on a un schéma en groupoïdes $(U, R, s, t, c)$
dont les morphismes $s$ et $t$ sont finis localement libres.
D'après Propriétés, Lemme \ref{properties-lemma-maximal-points-affine},
on peut trouver un ouvert affine contenant la réunion des orbites de
$u_1, \ldots, u_m$. Enfin, on peut appliquer
Groupoïdes, Lemme \ref{groupoids-lemma-find-invariant-affine}
pour conclure.
\end{proof}

\noindent
Le lemme suivant est un cas particulier du
Lemme \ref{lemma-invariant-affine-open-around-generic-point},
mais nous reprenons l'argument, car il est légèrement plus simple dans ce cas
(il évite d'utiliser le
Lemme \ref{lemma-finite-flat-over-almost-dense-subscheme-addendum}).

\begin{lemma}
\label{lemma-invariant-affine-open-around-generic-point-Noetherian}
Soit $(U, R, s, t, c)$ un schéma en groupoïdes sur un schéma $S$.
Supposons $s, t$ finis et $U$ localement noethérien, et soient $u_1, \ldots, u_m \in U$
des points dont les orbites sont constituées de points génériques de
composantes irréductibles de $U$. Il existe alors des sous-schémas invariants par $R$
$V' \subset V \subset U$ tels que
\begin{enumerate}
\item $u_1, \ldots, u_m \in V'$,
\item $V$ soit ouvert dans $U$,
\item $V'$ et $V$ soient affines,
\item $V' \subset V$ soit un épaississement,
\item les morphismes $s', t'$ de la restriction $(V', R', s', t', c')$
soient finis localement libres.
\end{enumerate}
\end{lemma}

\begin{proof}
Soit $\{u_{j1}, \ldots, u_{jn_j}\}$ l'orbite de $u_j$.
Soient $W' \subset W \subset U$ comme dans le
Lemme \ref{lemma-finite-flat-over-almost-dense-subscheme}.
Comme $U = t(s^{-1}(\overline{W}))$, on voit qu'il existe au moins
un $u_{ji} \in \overline{W}$. Comme $u_{ji}$ est un point générique
d'une composante irréductible et que $U$ est localement noethérien,
cela implique que $u_{ji} \in W$. Comme $W$ est invariant par $R$, on
conclut que $u_j \in W$ et, en fait, que l'orbite entière est contenue dans $W$.
D'après Cohomologie des schémas, Lemme
\ref{coherent-lemma-image-affine-finite-morphism-affine-Noetherian},
il suffit de trouver un sous-schéma ouvert affine $V'$ invariant par $R$
de $W'$ contenant $u_1, \ldots, u_m$ (car on peut alors prendre pour $V \subset W$
le sous-schéma ouvert correspondant, qui sera affine).
On peut donc remplacer $(U, R, s, t, c)$
par la restriction $(W', R', s', t', c')$ à $W'$.
Autrement dit, on peut supposer que l'on a un schéma en groupoïdes $(U, R, s, t, c)$
dont les morphismes $s$ et $t$ sont finis localement libres.
D'après Propriétés, Lemme \ref{properties-lemma-maximal-points-affine},
on peut trouver un ouvert affine contenant $\{u_{ij}\}$
(un schéma localement noethérien est quasi-séparé d'après
Propriétés, Lemme \ref{properties-lemma-locally-Noetherian-quasi-separated}).
Enfin, on peut appliquer
Groupoïdes, Lemme \ref{groupoids-lemma-find-invariant-affine}
pour conclure.
\end{proof}

\begin{lemma}
\label{lemma-find-affine-integral}
Soit $(U, R, s, t, c)$ un schéma en groupoïdes sur un schéma $S$
avec $s, t$ intégraux. Soit $g : U' \to U$ un morphisme intégral
tel que toute orbite sous $R$ dans $U$ rencontre $g(U')$. Soit $(U', R', s', t', c')$
la restriction de $R$ à $U'$. Si $u' \in U'$ appartient à un
ouvert affine invariant par $R'$, alors l'image $u \in U$ appartient
à un ouvert affine de $U$ invariant par $R$.
\end{lemma}

\begin{proof}
Soit $W' \subset U'$ un ouvert affine invariant par $R'$.
Posons $\tilde R = U' \times_{g, U, t} R$, muni des applications
$\text{pr}_0 : \tilde R \to U'$ et $h = s \circ \text{pr}_1 : \tilde R \to U$.
Observons que $\text{pr}_0$ et $h$ sont intégraux.
Il s'ensuit que $\tilde W = \text{pr}_0^{-1}(W')$ est affine.
Comme $W'$ est invariant par $R'$, l'image
$W = h(\tilde W)$ est ensemblistement invariante par $R$ et
$\tilde W = h^{-1}(W)$ ensemblistement (détails omis).
Ainsi, si l'on peut montrer que $W$ est ouvert, alors $W$ est un schéma
et le morphisme $\tilde W \to W$ est intégral surjectif,
ce qui implique que $W$ est affine d'après
Limites, Proposition \ref{limits-proposition-affine}.
Cependant, l'hypothèse selon laquelle les orbites rencontrent $U'$ implique
que $h : \tilde R \to U$ est surjectif. Comme un
morphisme intégral surjectif est submersif
(Topologie, Lemme \ref{topology-lemma-closed-morphism-quotient-topology}
et Morphismes, Lemme \ref{morphisms-lemma-integral-universally-closed}),
il s'ensuit que $W$ est ouvert.
\end{proof}

\noindent
Le lemme technique suivant produit des fonctions « presque » invariantes
dans la situation d'un groupoïde fini sur un
schéma quasi-affine.

\begin{lemma}
\label{lemma-find-almost-invariant-function}
Soit $(U, R, s, t, c)$ un schéma en groupoïdes avec $s, t$ finis et de
présentation finie. Soient $u_1, \ldots, u_m \in U$ des points dont les orbites sous $R$
sont constituées de points génériques de composantes irréductibles de $U$.
Soit $j : U \to \Spec(A)$ une immersion.
Soit $I \subset A$ un idéal tel que $j(U) \cap V(I) = \emptyset$
et que $V(I) \cup j(U)$ soit fermé dans $\Spec(A)$.
Il existe alors un élément $h \in I$ tel que $j^{-1}D(h)$
soit un sous-schéma ouvert affine de $U$ invariant par $R$ et contenant
$u_1, \ldots, u_m$.
\end{lemma}

\begin{proof}
Soient $u_1, \ldots, u_m \in V' \subset V \subset U$ comme dans le
Lemme \ref{lemma-invariant-affine-open-around-generic-point}.
Comme $U \setminus V$ est fermé dans $U$, que $j$ est une immersion et que $V(I) \cup j(U)$
est fermé dans $\Spec(A)$, on peut trouver un idéal
$J \subset I$ tel que $V(J) = V(I) \cup j(U \setminus V)$.
On peut par exemple prendre l'idéal des éléments de $I$ qui
s'annulent sur $j(U \setminus V)$. On peut donc remplacer
$(U, R, s, t, c)$, $j : U \to \Spec(A)$ et $I$ par
$(V', R', s', t', c')$, $j|_{V'} : V' \to \Spec(A)$ et $J$.
Autrement dit, on peut supposer que $U$ est affine et que
$s$ et $t$ sont finis localement libres.
Prenons un quelconque $f \in I$ qui ne s'annule en aucun des
points des orbites sous $R$ de $u_1, \ldots, u_m$
(Algèbre, Lemme \ref{algebra-lemma-silly}). Considérons
$$
g = \text{Norm}_s(t^\sharp(j^\sharp(f))) \in \Gamma(U, \mathcal{O}_U)
$$
Comme $f \in I$ et que $V(I) \cup j(U)$ est fermé, on voit que
$U \cap D(f) \to D(f)$ est une immersion fermée. Ainsi, $f^ng$ est
l'image d'un élément $h \in I$ pour un certain $n > 0$. Nous affirmons que $h$ convient.
En effet, nous avons vu dans
Groupoïdes, Lemme \ref{groupoids-lemma-determinant-trick},
que $g$ est une fonction invariante par $R$ ; ainsi, $D(g) \subset U$
est invariant par $R$. Comme $f$ ne s'annule pas sur l'orbite de $u_j$,
la fonction $g$ ne s'annule pas en $u_j$. De plus, on a
$V(g) \supset V(j^\sharp(f))$ et donc $j^{-1}D(h) = D(g)$.
\end{proof}

\begin{lemma}
\label{lemma-no-specializations-map-to-same-point}
Soit $(U, R, s, t, c)$ un schéma en groupoïdes. Si $s, t$ sont finis,
et si $u, u' \in R$ sont des points distincts de la même orbite,
alors $u'$ n'est pas une spécialisation de $u$.
\end{lemma}

\begin{proof}
Soit $r \in R$ tel que $s(r) = u$ et $t(r) = u'$. Si $u \leadsto u'$
alors nous pouvons trouver une spécialisation non triviale $r \leadsto r'$ telle que
$s(r') = u'$, voir 
Schémas, lemme \ref{schemes-lemma-quasi-compact-closed}.
Posons $u'' = t(r')$. Remarquons que $u'' \not = u'$ puisqu'il n'existe pas de
spécialisations dans les fibres d'un morphisme fini.
Nous pouvons donc continuer et trouver une spécialisation non triviale
$r' \leadsto r''$ telle que $s(r'') = u''$, et ainsi de suite. Cela montre que
l'orbite de $u$ contient une suite infinie
$u \leadsto u' \leadsto u'' \leadsto \ldots$
de spécialisations, ce qui est absurde puisque l'orbite
$t(s^{-1}(\{u\}))$ est finie.
\end{proof}

\begin{lemma}
\label{lemma-get-affine}
Soit $j : V \to \Spec(A)$ une immersion quasi-compacte de schémas.
Soit $f \in A$ tel que $j^{-1}D(f)$ soit affine et que $j(V) \cap V(f)$
soit fermé. Alors $V$ est affine.
\end{lemma}

\begin{proof}
Cela résulte de Morphismes, lemme \ref{morphisms-lemma-get-affine},
mais nous allons aussi en donner une démonstration directe.
Posons $A' = \Gamma(V, \mathcal{O}_V)$. Alors $j' : V \to \Spec(A')$ est une
immersion ouverte quasi-compacte, voir
Propriétés, lemme \ref{properties-lemma-quasi-affine}.
Soit $f' \in A'$ l'image de $f$. Alors $(j')^{-1}D(f') = j^{-1}D(f)$
est affine. D'autre part, $j'(V) \cap V(f')$ est un sous-schéma de
$\Spec(A')$ qui s'envoie isomorphiquement sur le sous-schéma fermé
$j(V) \cap V(f)$ de $\Spec(A)$. Il est donc fermé dans $\Spec(A')$,
par exemple d'après Schémas, lemme \ref{schemes-lemma-section-immersion}.
Ainsi, nous pouvons remplacer $A$ par $A'$ et supposer que $j$ est une immersion ouverte
et que $A = \Gamma(V, \mathcal{O}_V)$.

\medskip\noindent
Dans ce cas, nous affirmons que $j(V) = \Spec(A)$, ce qui achève la démonstration.
Sinon, nous pouvons trouver un ouvert affine principal $D(g) \subset \Spec(A)$
qui rencontre le complémentaire et évite la partie fermée $j(V) \cap V(f)$.
Remarquons que $j$ envoie $j^{-1}D(f)$ isomorphiquement sur $D(f)$, voir
Propriétés, lemme \ref{properties-lemma-invert-f-affine}.
Ainsi $D(g)$ rencontre $V(f)$. D'autre part, $j^{-1}D(g)$
est un ouvert principal de l'ouvert affine $j^{-1}D(f)$, donc est affine.
Par conséquent, en appliquant de nouveau
Propriétés, lemme \ref{properties-lemma-invert-f-affine},
nous voyons que $D(g)$ est isomorphe à $j^{-1}D(g) \subset j^{-1}D(f)$,
ce qui implique que $D(g) \subset D(f)$. Cette contradiction achève
la démonstration.
\end{proof}

\begin{lemma}
\label{lemma-find-affine-codimension-1}
Soit $(U, R, s, t, c)$ un schéma en groupoïdes. Soit $u \in U$. Supposons que
\begin{enumerate}
\item $s, t$ soient des morphismes finis,
\item $U$ soit séparé et localement noethérien,
\item $\dim(\mathcal{O}_{U, u'}) \leq 1$ pour tout point $u'$
de l'orbite de $u$.
\end{enumerate}
Alors $u$ est contenu dans un ouvert affine de $U$ invariant par $R$.
\end{lemma}

\begin{proof}
L'orbite sous $R$ de $u$ est finie. D'après les conditions (2) et (3), elle est contenue
dans un ouvert affine $U'$ de $U$, voir
Variétés, proposition
\ref{varieties-proposition-finite-set-of-points-of-codim-1-in-affine}.
Alors $t(s^{-1}(U \setminus U'))$ est une partie fermée et invariante par $R$
de $U$ qui ne contient pas $u$. Ainsi
$U \setminus t(s^{-1}(U \setminus U'))$ est un ouvert invariant par $R$
de $U'$ contenant $u$.
En remplaçant $U$ par cet ouvert, nous pouvons supposer que $U$ est quasi-affine.

\medskip\noindent
D'après le lemme \ref{lemma-find-affine-integral}, nous pouvons remplacer $U$ par son réduit
et supposer que $U$ est réduit. Cela signifie que les sous-schémas invariants par $R$
$W' \subset W \subset U$ du
lemme \ref{lemma-finite-flat-over-almost-dense-subscheme}
sont égaux : $W' = W$. Comme $U = t(s^{-1}(\overline{W}))$, un point
$u'$ appartenant à l'orbite sous $R$ de $u$ est contenu dans $\overline{W}$
et, d'après le lemme \ref{lemma-find-affine-integral},
nous pouvons remplacer $U$ par $\overline{W}$ et $u$ par $u'$.
Nous pouvons donc supposer qu'il existe
un sous-schéma ouvert dense $W \subset U$ invariant par $R$ tel que
les morphismes $s_W, t_W$ de la restriction $(W, R_W, s_W, t_W, c_W)$ soient
finis localement libres.

\medskip\noindent
Si $u \in W$, nous avons terminé d'après
Groupoïdes, lemme \ref{groupoids-lemma-find-invariant-affine}
(car $W$ est quasi-affine, donc tout ensemble fini de points
de $W$ est contenu dans un ouvert affine, voir
Propriétés, lemme \ref{properties-lemma-ample-finite-set-in-affine}).
Nous supposons donc que $u \not \in W$ et, par suite, qu'aucun des points de
l'orbite de $u$ n'appartient à $W$. Soit $\xi \in U$
un point qui se spécialise non trivialement en un point $u'$ de l'orbite
de $u$. Puisqu'il n'y a pas de spécialisations entre les points de
l'orbite de $u$ (lemme \ref{lemma-no-specializations-map-to-same-point}),
nous voyons que $\xi$ n'appartient pas à l'orbite.
D'après l'hypothèse (3), nous voyons que $\xi$ est un point générique de $U$
et donc que $\xi \in W$.
Comme $U$ est noethérien, il n'y a qu'un nombre fini de ces
points $\xi_1, \ldots, \xi_m \in W$. Puisque $s_W, t_W$ sont plats, l'orbite
de chaque $\xi_j$ est formée de points génériques de composantes irréductibles
de $W$ (et donc de $U$).

\medskip\noindent
Soit $j : U \to \Spec(A)$ une immersion de $U$ dans un schéma affine
(cela est possible puisque $U$ est quasi-affine). Soit $J \subset A$
un idéal tel que $V(J) \cap j(W) = \emptyset$ et que $V(J) \cup j(W)$
soit fermé. Appliquons le lemme \ref{lemma-find-almost-invariant-function}
au schéma en groupoïdes $(W, R_W, s_W, t_W, c_W)$, au morphisme
$j|_W : W \to \Spec(A)$, aux points $\xi_j$ et à l'idéal $J$
pour trouver $f \in J$ tel que $(j|_W)^{-1}D(f)$ soit un ouvert invariant par $R_W$
affine contenant $\xi_j$ pour tout $j$. Puisque $f \in J$,
nous voyons que $j^{-1}D(f) \subset W$, c'est-à-dire que $j^{-1}D(f)$ est
un ouvert affine de $U$ invariant par $R$ et contenu dans $W$
et contenant tous les $\xi_j$.

\medskip\noindent
Soit $Z$ le sous-schéma fermé réduit induit sur
$$
U \setminus j^{-1}D(f) = j^{-1}V(f).
$$
Alors $Z$ est, du point de vue ensembliste,
invariant par $R$ (mais il peut ne pas être invariant par $R$ au sens schématique).
Soit $(Z, R_Z, s_Z, t_Z, c_Z)$ la restriction de $R$ à $Z$.
Puisque $Z \to U$ est fini, il s'ensuit que $s_Z$ et $t_Z$ sont finis.
Puisque $u \in Z$, l'orbite de $u$ est contenue dans $Z$ et coïncide avec
l'orbite sous $R_Z$ de $u$, considéré comme un point de $Z$. Puisque
$\dim(\mathcal{O}_{U, u'}) \leq 1$ et puisque $\xi_j \not \in Z$
pour tout $j$, nous voyons que $\dim(\mathcal{O}_{Z, u'}) \leq 0$ pour
tout $u'$ de l'orbite de $u$. Autrement dit, l'orbite sous $R_Z$ de $u$
est formée de points génériques de composantes irréductibles de $Z$.

\medskip\noindent
Soit $I \subset A$ un idéal tel que $V(I) \cap j(U) =\emptyset$
et que $V(I) \cup j(U)$ soit fermé. Appliquons
le lemme \ref{lemma-find-almost-invariant-function} au
schéma en groupoïdes $(Z, R_Z, s_Z, t_Z, c_Z)$, à la restriction $j|_Z$,
à l'idéal $I$ et au point $u \in Z$, pour obtenir $h \in I$ tel que
$j^{-1}D(h) \cap Z$ soit un ouvert affine invariant par $R_Z$ et contenant $u$.

\medskip\noindent
Considérons la fonction invariante par $R_W$ (Groupoïdes, lemme
\ref{groupoids-lemma-determinant-trick})
$$
g = 
\text{Norme}_{s_W}(t_W^\sharp(j^\sharp(h)|_W)) \in \Gamma(W, \mathcal{O}_W)
$$
(Dans ce qui suit, nous n'avons besoin que de la restriction de $g$ à $j^{-1}D(f)$ et,
dans ce cas, la norme est prise le long d'un morphisme fini localement libre entre schémas affines.)
Nous affirmons que
$$
V = (W_g \cap j^{-1}D(f)) \cup (j^{-1}D(h) \cap Z)
$$
est un ouvert affine de $U$ invariant par $R$, ce qui achève la démonstration du lemme.
Il est invariant par $R$ du point de vue ensembliste par construction. Comme $V$ est une
partie constructible, pour voir qu'il est ouvert il suffit de montrer qu'il est
stable par générisation dans $U$ (Topologie, lemme
\ref{topology-lemma-characterize-closed-Noetherian}
ou, plus généralement,
Topologie, lemme
\ref{topology-lemma-constructible-stable-specialization-closed}).
Puisque $W_g \cap j^{-1}D(f)$ est ouvert dans $U$, il suffit de considérer
une spécialisation $u_1 \leadsto u_2$ de $U$ telle que
$u_2 \in j^{-1}D(h) \cap Z$.
Cela signifie que $h$ est non nul en $j(u_2)$ et que $u_2 \in Z$.
Si $u_1 \in Z$, alors $j(u_1) \leadsto j(u_2)$ et, puisque
$h$ est non nul en $j(u_2)$, il est non nul en $j(u_1)$, ce qui
implique que $u_1 \in V$. Si $u_1 \not \in Z$ et n'appartient pas non plus
à $W_g \cap j^{-1}D(f)$, alors $u_1 \in W$, $u_1 \not \in W_g$
car le complémentaire de $Z = j^{-1}V(f)$ est contenu dans $W \cap j^{-1}D(f)$.
Il existe donc un point $r_1 \in R$ tel que $s(r_1) = u_1$
et tel que $h$ soit nul en $t(r_1)$. Puisque $s$ est fini, nous
pouvons trouver une spécialisation $r_1 \leadsto r_2$ telle que $s(r_2) = u_2$.
Mais nous en déduisons alors que $h$ est nul en $u'_2 = t(r_2)$,
ce qui contredit le fait que $j^{-1}D(h) \cap Z$ est
invariant par $R$ et que $u_2$ lui appartient. Ainsi $V$ est ouvert.

\medskip\noindent
Observons que $V \subset j^{-1}D(h)$ pour notre fonction $h \in I$.
Nous obtenons donc une immersion
$$
j' : V \longrightarrow \Spec(A_h)
$$
Soit $f' \in A_h$ l'image de $f$. Alors $(j')^{-1}D(f')$
est l'ouvert principal déterminé par $g$ dans l'ouvert
affine $j^{-1}D(f)$ de $U$.
Ainsi $(j')^{-1}D(f)$ est affine. Enfin,
$j'(V) \cap V(f') = j'(j^{-1}D(h) \cap Z)$
est fermé dans $\Spec(A_h/(f')) = \Spec((A/f)_h) = D(h) \cap V(f)$
en vertu du choix de $h \in I$ et de l'idéal $I$. Nous pouvons donc appliquer le
lemme \ref{lemma-get-affine}
pour conclure que $V$ est affine, comme affirmé ci-dessus.
\end{proof}







\section{Descente des morphismes ind-quasi-affines}
\label{section-ind-quasi-affine}

\noindent
Les morphismes ind-quasi-affines ont été définis dans
Compléments sur les morphismes, section \ref{more-morphisms-section-ind-quasi-affine}.
Cette section est l'analogue de
Descente, section \ref{descent-section-quasi-affine},
pour les morphismes ind-quasi-affines.

\medskip\noindent
Soit $X$ un schéma quasi-séparé. Soit $E \subset X$ une partie
qui est l'intersection d'une famille non vide d'ouverts quasi-compacts de $X$.
Écrivons $E = \bigcap_{i \in I} U_i$, où $U_i \subset X$ est un ouvert quasi-compact
et où $I$ est non vide.
En ajoutant les intersections finies, nous pouvons supposer que, pour $i, j \in I$,
il existe $k \in I$ tel que $U_k \subset U_i \cap U_j$.
Dans cette situation, nous avons
\begin{equation}
\label{equation-sections-of-intersection}
\Gamma(E, \mathcal{F}|_E) = \colim \Gamma(U_i, \mathcal{F}|_{U_i})
\end{equation}
pour tout faisceau $\mathcal{F}$ défini sur $X$. En effet, fixons $i_0 \in I$
et remplaçons $X$ par $U_{i_0}$ et $I$ par
$\{i \in I \mid U_i \subset U_{i_0}\}$. Alors $X$ est quasi-compact
et quasi-séparé, donc est un espace spectral, voir
Propriétés, lemme
\ref{properties-lemma-quasi-compact-quasi-separated-spectral}.
Nous voyons alors que l'égalité résulte de
Topologie, lemme \ref{topology-lemma-make-spectral-space}, et de
Faisceaux, lemme \ref{sheaves-lemma-descend-opens}.
(En fait, la formule vaut aussi pour les groupes de cohomologie supérieurs
lorsque $\mathcal{F}$ est abélien, voir
Cohomologie, lemme \ref{cohomology-lemma-colimit}.)

\begin{lemma}
\label{lemma-sits-in-functions}
Soit $X$ un schéma ind-quasi-affine. Soit $E \subset X$ une
intersection d'une famille non vide d'ouverts quasi-compacts de $X$.
Posons $A = \Gamma(E, \mathcal{O}_X|_E)$ et $Y = \Spec(A)$.
Alors le morphisme canonique
$$
j : (E, \mathcal{O}_X|_E) \longrightarrow (Y, \mathcal{O}_Y)
$$
fourni par Schémas, lemme \ref{schemes-lemma-morphism-into-affine},
détermine un isomorphisme
$(E, \mathcal{O}_X|_E) \to (E', \mathcal{O}_Y|_{E'})$
où $E' \subset Y$ est une intersection d'ouverts quasi-compacts.
Si $W \subset E$ est ouvert dans $X$, alors $j(W)$ est ouvert dans $Y$.
\end{lemma}

\begin{proof}
Remarquons que $(E, \mathcal{O}_X|_E)$ est un espace localement annelé, de sorte que
Schémas, lemme \ref{schemes-lemma-morphism-into-affine}, s'applique
à $A \to \Gamma(E, \mathcal{O}_X|_E)$. Écrivons $E = \bigcap_{i \in I} U_i$,
avec $I \not = \emptyset$ et $U_i \subset X$ ouvert quasi-compact.
Nous pouvons supposer, et nous le faisons, que pour $i, j \in I$ il existe $k \in I$ tel que
$U_k \subset U_i \cap U_j$. Posons $A_i = \Gamma(U_i, \mathcal{O}_{U_i})$.
Nous obtenons les diagrammes commutatifs
$$
\xymatrix{
(E, \mathcal{O}_X|_E) \ar[r] \ar[d] &
(\Spec(A), \mathcal{O}_{\Spec(A)}) \ar[d] \\
(U_i, \mathcal{O}_{U_i}) \ar[r] &
(\Spec(A_i), \mathcal{O}_{\Spec(A_i)})
}
$$
Puisque $U_i$ est quasi-affine, nous voyons que $U_i \to \Spec(A_i)$
est une immersion ouverte quasi-compacte. D'autre part,
$A = \colim A_i$. Ainsi $\Spec(A) = \lim \Spec(A_i)$ comme espaces
topologiques (Limites, lemme \ref{limits-lemma-topology-limit}). Puisque
$E = \lim U_i$ (d'après Topologie, lemme \ref{topology-lemma-make-spectral-space}),
nous voyons que $E \to \Spec(A)$ est un homéomorphisme sur son
image $E'$ et que $E'$ est l'intersection des images réciproques
des ouverts $U_i \subset \Spec(A_i)$ dans $\Spec(A)$. Pour tout
$e \in E$, l'anneau local $\mathcal{O}_{X, e}$ est
$\mathcal{O}_{U_i, e}$, qui coïncide avec l'anneau local correspondant sur $\Spec(A)$.

\medskip\noindent
Pour démontrer la dernière assertion du lemme, raisonnons comme suit.
Choisissons $i, j \in I$ tels que $U_i \subset U_j$.
Considérons les diagrammes commutatifs suivants
$$
\xymatrix{
U_i \ar[r] \ar[d] & \Spec(A_i) \ar[d] \\
U_i \ar[r] & \Spec(A_j)
}
\quad\quad
\xymatrix{
W \ar[r] \ar[d] & \Spec(A_i) \ar[d] \\
W \ar[r] & \Spec(A_j)
}
\quad\quad
\xymatrix{
W \ar[r] \ar[d] & \Spec(A) \ar[d] \\
W \ar[r] & \Spec(A_j)
}
$$
D'après Propriétés, lemme
\ref{properties-lemma-cartesian-diagram-quasi-affine},
le premier diagramme est cartésien. Le deuxième l'est donc aussi.
En passant à la limite, nous trouvons que le troisième diagramme
est cartésien, si bien que la flèche horizontale supérieure de ce diagramme est une immersion ouverte.
\end{proof}

\begin{lemma}
\label{lemma-affine-base-change}
Supposons donné un diagramme cartésien
$$
\xymatrix{
X \ar[d]_f \ar[r] & \Spec(B) \ar[d] \\
Y \ar[r] & \Spec(A)
}
$$
de schémas. Soit $E \subset Y$ une intersection d'une famille non vide
d'ouverts quasi-compacts de $Y$. Alors
$$
\Gamma(f^{-1}(E), \mathcal{O}_X|_{f^{-1}(E)}) =
\Gamma(E, \mathcal{O}_Y|_E) \otimes_A B
$$
pourvu que $Y$ soit quasi-séparé et que $A \to B$ soit plat.
\end{lemma}

\begin{proof}
Écrivons $E = \bigcap_{i \in I} V_i$, où $V_i \subset Y$ est un ouvert quasi-compact.
Nous pouvons supposer, et nous le faisons, que pour $i, j \in I$ il existe $k \in I$ tel que
$V_k \subset V_i \cap V_j$. Nous avons alors de même
$f^{-1}(E) = \bigcap_{i \in I} f^{-1}(V_i)$ dans $X$.
Le résultat découle donc de l'équation (\ref{equation-sections-of-intersection})
et du résultat correspondant pour $V_i$ et $f^{-1}(V_i)$, qui est
Cohomologie des schémas, lemme \ref{coherent-lemma-flat-base-change-cohomology}.
\end{proof}

\begin{lemma}[Gabber]
\label{lemma-ind-quasi-affine}
Soit $S$ un schéma. Soit $\{X_i \to S\}_{i\in I}$ un recouvrement fpqc.
Soit $(V_i/X_i, \varphi_{ij})$ une donnée de descente relativement à
$\{X_i \to S\}$, voir Descente, définition
\ref{descent-definition-descent-datum-for-family-of-morphisms}. 
Si chaque morphisme $V_i \to X_i$ est ind-quasi-affine, alors la donnée de descente
est effective.
\end{lemma}

\begin{proof}
Être ind-quasi-affine est une propriété des morphismes de schémas
qui est préservée par tout changement de base, voir
Compléments sur les morphismes, lemme
\ref{more-morphisms-lemma-base-change-ind-quasi-affine}.
Ainsi Descente, lemme \ref{descent-lemma-descending-types-morphisms}, s'applique,
et il suffit de démontrer l'énoncé du lemme
dans le cas où le recouvrement fpqc est constitué du singleton
$\{X \to S\}$, formé d'un morphisme plat surjectif de schémas affines.
Écrivons $X = \Spec(A)$ et $S = \Spec(R)$, de sorte
que $R \to A$ est un homomorphisme d'anneaux fidèlement plat.
Soit $(V, \varphi)$ une donnée de descente relative à $X$ au-dessus de $S$
et supposons que $V \to X$ soit ind-quasi-affine, autrement dit que
$V$ soit ind-quasi-affine.

\medskip\noindent
Soit $(U, R, s, t, c)$ le schéma en groupoïdes au-dessus de $S$ où
$U = X$ et $R = X \times_S X$, et où $s$, $t$, $c$ sont les morphismes usuels.
D'après Groupoïdes, lemme \ref{groupoids-lemma-cartesian-equivalent-descent-datum},
la paire $(V, \varphi)$ correspond à un morphisme cartésien
$(U', R', s', t', c') \to (U, R, s, t, c)$ de schémas en groupoïdes.
Soit $u' \in U'$ un point quelconque. D'après
Groupoïdes, lemmes \ref{groupoids-lemma-constructing-invariant-opens},
\ref{groupoids-lemma-first-observation} et
\ref{groupoids-lemma-second-observation},
nous pouvons choisir $u' \in W \subset E \subset U'$
où $W$ est ouvert et invariant par $R'$, et où
$E$ est invariant par $R'$ du point de vue ensembliste et
est une intersection d'une famille non vide d'ouverts quasi-compacts.

\medskip\noindent
En revenant à $(V, \varphi)$, pour tout $v \in V$ nous pouvons choisir
$v \in W \subset E \subset V$ de sorte que les propriétés suivantes soient vérifiées :
(a) $W$ est ouvert et $\varphi(W \times_S X) = X \times_S W$, et
(b) $E$ est une intersection d'ouverts quasi-compacts et
$\varphi(E \times_S X) = X \times_S E$ du point de vue ensembliste.
Ici, la notation $E \times_S X$ désigne
l'image réciproque de $E$ dans $V \times_S X$ par le morphisme de projection, et
de même pour $X \times_S E$. D'après le lemme \ref{lemma-affine-base-change},
cela implique que $\varphi$ définit un isomorphisme
\begin{align*}
\Gamma(E, \mathcal{O}_V|_E) \otimes_R A
& =
\Gamma(E \times_S X, \mathcal{O}_{V \times_S X}|_{E \times_S X}) \\
& \to
\Gamma(X \times_S E, \mathcal{O}_{X \times_S V}|_{X \times_S E}) \\
& =
A \otimes_R \Gamma(E, \mathcal{O}_V|_E)
\end{align*}
d'algèbres sur $A \otimes_R A$, que nous appellerons $\psi$.
La condition de cocycle pour $\varphi$
se traduit par la condition de cocycle pour $\psi$ comme dans
Descente, définition \ref{descent-definition-descent-datum-modules}
(détails omis). D'après Descente, proposition
\ref{descent-proposition-descent-module},
nous trouvons une $R$-algèbre $R'$ et un isomorphisme
$\chi : R' \otimes_R A \to \Gamma(E, \mathcal{O}_V|_E)$
d'$A$-algèbres, compatible avec $\psi$ et avec la
donnée de descente canonique sur $R' \otimes_R A$.

\medskip\noindent
D'après le lemme \ref{lemma-sits-in-functions}, nous obtenons un ``plongement'' canonique
$$
j : (E, \mathcal{O}_V|_E) \longrightarrow
\Spec(\Gamma(E, \mathcal{O}_V|_E)) = \Spec(R' \otimes_R A)
$$
d'espaces localement annelés. La construction de cette application est canonique
et nous obtenons un diagramme commutatif
$$
\xymatrix{
& E \times_S X \ar[rr]_\varphi \ar[ld] \ar[rd]^{j'} & &
X \times_S E \ar[rd] \ar[ld]_{j''} \\
E \ar[rd]^j  & &
\Spec(R' \otimes_R A \otimes_R A) \ar[ld] \ar[rd] & &
E \ar[ld]_j \\
& \Spec(R' \otimes_R A) \ar[rd] && \Spec(R' \otimes_R A) \ar[ld] \\
& & \Spec(R')
}
$$
où $j'$ et $j''$ proviennent de la même construction appliquée à
$E \times_S X \subset V \times_S X$ et à $X \times_S E \subset X \times_S V$
au moyen de $\chi$ et des identifications utilisées pour construire $\psi$.
Il s'ensuit que $j(W)$ est un sous-schéma ouvert de $\Spec(R' \otimes_R A)$
dont les images réciproques par les deux projections
$\Spec(R' \otimes_R A \otimes_R A) \to \Spec(R' \otimes_R A)$
sont égales. D'après Descente, lemme \ref{descent-lemma-open-fpqc-covering},
nous trouvons un ouvert $W_0 \subset \Spec(R')$ dont le changement de base
à $\Spec(A)$ est $j(W)$. En examinant le diagramme ci-dessus,
nous voyons que la donnée de descente $(W, \varphi|_{W \times_S X})$
est effective. D'après Descente, lemme
\ref{descent-lemma-effective-for-fpqc-is-local-upstairs},
nous voyons que notre donnée de descente est effective.
\end{proof}






\begin{multicols}{2}[\section{Autres chapitres}]
\noindent
Préliminaires
\begin{enumerate}
\item \hyperref[introduction-section-phantom]{Introduction}
\item \hyperref[conventions-section-phantom]{Conventions}
\item \hyperref[sets-section-phantom]{Théorie des ensembles}
\item \hyperref[categories-section-phantom]{Catégories}
\item \hyperref[topology-section-phantom]{Topologie}
\item \hyperref[sheaves-section-phantom]{Faisceaux sur les espaces}
\item \hyperref[sites-section-phantom]{Sites et faisceaux}
\item \hyperref[stacks-section-phantom]{Champs}
\item \hyperref[fields-section-phantom]{Corps}
\item \hyperref[algebra-section-phantom]{Algèbre commutative}
\item \hyperref[brauer-section-phantom]{Groupes de Brauer}
\item \hyperref[homology-section-phantom]{Algèbre homologique}
\item \hyperref[derived-section-phantom]{Catégories dérivées}
\item \hyperref[simplicial-section-phantom]{Méthodes simpliciales}
\item \hyperref[more-algebra-section-phantom]{Compléments d'algèbre}
\item \hyperref[smoothing-section-phantom]{Lissage des morphismes d'anneaux}
\item \hyperref[modules-section-phantom]{Faisceaux de modules}
\item \hyperref[sites-modules-section-phantom]{Modules sur les sites}
\item \hyperref[injectives-section-phantom]{Objets injectifs}
\item \hyperref[cohomology-section-phantom]{Cohomologie des faisceaux}
\item \hyperref[sites-cohomology-section-phantom]{Cohomologie sur les sites}
\item \hyperref[dga-section-phantom]{Algèbre différentielle graduée}
\item \hyperref[dpa-section-phantom]{Algèbre à puissances divisées}
\item \hyperref[sdga-section-phantom]{Faisceaux différentiels gradués}
\item \hyperref[hypercovering-section-phantom]{Hyperrecouvrements}
\end{enumerate}
Schémas
\begin{enumerate}
\setcounter{enumi}{25}
\item \hyperref[schemes-section-phantom]{Schémas}
\item \hyperref[constructions-section-phantom]{Constructions de schémas}
\item \hyperref[properties-section-phantom]{Propriétés des schémas}
\item \hyperref[morphisms-section-phantom]{Morphismes de schémas}
\item \hyperref[coherent-section-phantom]{Cohomologie des schémas}
\item \hyperref[divisors-section-phantom]{Diviseurs}
\item \hyperref[limits-section-phantom]{Limites de schémas}
\item \hyperref[varieties-section-phantom]{Variétés}
\item \hyperref[topologies-section-phantom]{Topologies sur les schémas}
\item \hyperref[descent-section-phantom]{Descente}
\item \hyperref[perfect-section-phantom]{Catégories dérivées de schémas}
\item \hyperref[more-morphisms-section-phantom]{Compléments sur les morphismes}
\item \hyperref[flat-section-phantom]{Compléments sur la platitude}
\item \hyperref[groupoids-section-phantom]{Schémas en groupoïdes}
\item \hyperref[more-groupoids-section-phantom]{Compléments sur les schémas en groupoïdes}
\item \hyperref[etale-section-phantom]{Morphismes étales de schémas}
\end{enumerate}
Thèmes de théorie des schémas
\begin{enumerate}
\setcounter{enumi}{41}
\item \hyperref[chow-section-phantom]{Homologie de Chow}
\item \hyperref[intersection-section-phantom]{Théorie de l'intersection}
\item \hyperref[pic-section-phantom]{Schémas de Picard des courbes}
\item \hyperref[weil-section-phantom]{Théories cohomologiques de Weil}
\item \hyperref[adequate-section-phantom]{Modules adéquats}
\item \hyperref[dualizing-section-phantom]{Complexes dualisants}
\item \hyperref[duality-section-phantom]{Dualité pour les schémas}
\item \hyperref[discriminant-section-phantom]{Discriminants et différentes}
\item \hyperref[derham-section-phantom]{Cohomologie de de Rham}
\item \hyperref[local-cohomology-section-phantom]{Cohomologie locale}
\item \hyperref[algebraization-section-phantom]{Géométrie algébrique et formelle}
\item \hyperref[curves-section-phantom]{Courbes algébriques}
\item \hyperref[resolve-section-phantom]{Résolution des surfaces}
\item \hyperref[models-section-phantom]{Réduction semi-stable}
\item \hyperref[functors-section-phantom]{Foncteurs et morphismes}
\item \hyperref[equiv-section-phantom]{Catégories dérivées de variétés}
\item \hyperref[pione-section-phantom]{Groupes fondamentaux de schémas}
\item \hyperref[etale-cohomology-section-phantom]{Cohomologie étale}
\item \hyperref[crystalline-section-phantom]{Cohomologie cristalline}
\item \hyperref[proetale-section-phantom]{Cohomologie pro-étale}
\item \hyperref[relative-cycles-section-phantom]{Cycles relatifs}
\item \hyperref[more-etale-section-phantom]{Compléments de cohomologie étale}
\item \hyperref[trace-section-phantom]{La formule des traces}
\end{enumerate}
Espaces algébriques
\begin{enumerate}
\setcounter{enumi}{64}
\item \hyperref[spaces-section-phantom]{Espaces algébriques}
\item \hyperref[spaces-properties-section-phantom]{Propriétés des espaces algébriques}
\item \hyperref[spaces-morphisms-section-phantom]{Morphismes d'espaces algébriques}
\item \hyperref[decent-spaces-section-phantom]{Espaces algébriques décents}
\item \hyperref[spaces-cohomology-section-phantom]{Cohomologie des espaces algébriques}
\item \hyperref[spaces-limits-section-phantom]{Limites d'espaces algébriques}
\item \hyperref[spaces-divisors-section-phantom]{Diviseurs sur les espaces algébriques}
\item \hyperref[spaces-over-fields-section-phantom]{Espaces algébriques sur des corps}
\item \hyperref[spaces-topologies-section-phantom]{Topologies sur les espaces algébriques}
\item \hyperref[spaces-descent-section-phantom]{Descente et espaces algébriques}
\item \hyperref[spaces-perfect-section-phantom]{Catégories dérivées d'espaces}
\item \hyperref[spaces-more-morphisms-section-phantom]{Compléments sur les morphismes d'espaces}
\item \hyperref[spaces-flat-section-phantom]{Platitude sur les espaces algébriques}
\item \hyperref[spaces-groupoids-section-phantom]{Groupoïdes dans les espaces algébriques}
\item \hyperref[spaces-more-groupoids-section-phantom]{Compléments sur les groupoïdes dans les espaces}
\item \hyperref[bootstrap-section-phantom]{Amorçage}
\item \hyperref[spaces-pushouts-section-phantom]{Sommes amalgamées d'espaces algébriques}
\end{enumerate}
Thèmes de géométrie
\begin{enumerate}
\setcounter{enumi}{81}
\item \hyperref[spaces-chow-section-phantom]{Groupes de Chow des espaces}
\item \hyperref[groupoids-quotients-section-phantom]{Quotients de groupoïdes}
\item \hyperref[spaces-more-cohomology-section-phantom]{Compléments sur la cohomologie des espaces}
\item \hyperref[spaces-simplicial-section-phantom]{Espaces simpliciaux}
\item \hyperref[spaces-duality-section-phantom]{Dualité pour les espaces}
\item \hyperref[formal-spaces-section-phantom]{Espaces algébriques formels}
\item \hyperref[restricted-section-phantom]{Algébrisation des espaces formels}
\item \hyperref[spaces-resolve-section-phantom]{Résolution des surfaces revisitée}
\end{enumerate}
Théorie des déformations
\begin{enumerate}
\setcounter{enumi}{89}
\item \hyperref[formal-defos-section-phantom]{Théorie formelle des déformations}
\item \hyperref[defos-section-phantom]{Théorie des déformations}
\item \hyperref[cotangent-section-phantom]{Le complexe cotangent}
\item \hyperref[examples-defos-section-phantom]{Problèmes de déformation}
\end{enumerate}
Champs algébriques
\begin{enumerate}
\setcounter{enumi}{93}
\item \hyperref[algebraic-section-phantom]{Champs algébriques}
\item \hyperref[examples-stacks-section-phantom]{Exemples de champs}
\item \hyperref[stacks-sheaves-section-phantom]{Faisceaux sur les champs algébriques}
\item \hyperref[criteria-section-phantom]{Critères de représentabilité}
\item \hyperref[artin-section-phantom]{Axiomes d'Artin}
\item \hyperref[quot-section-phantom]{Espaces de Quot et de Hilbert}
\item \hyperref[stacks-properties-section-phantom]{Propriétés des champs algébriques}
\item \hyperref[stacks-morphisms-section-phantom]{Morphismes de champs algébriques}
\item \hyperref[stacks-limits-section-phantom]{Limites de champs algébriques}
\item \hyperref[stacks-cohomology-section-phantom]{Cohomologie des champs algébriques}
\item \hyperref[stacks-perfect-section-phantom]{Catégories dérivées de champs}
\item \hyperref[stacks-introduction-section-phantom]{Introduction aux champs algébriques}
\item \hyperref[stacks-more-morphisms-section-phantom]{Compléments sur les morphismes de champs}
\item \hyperref[stacks-geometry-section-phantom]{La géométrie des champs}
\end{enumerate}
Thèmes de théorie des espaces de modules
\begin{enumerate}
\setcounter{enumi}{107}
\item \hyperref[moduli-section-phantom]{Champs de modules}
\item \hyperref[moduli-curves-section-phantom]{Espaces de modules de courbes}
\end{enumerate}
Divers
\begin{enumerate}
\setcounter{enumi}{109}
\item \hyperref[examples-section-phantom]{Exemples}
\item \hyperref[exercises-section-phantom]{Exercices}
\item \hyperref[guide-section-phantom]{Guide bibliographique}
\item \hyperref[desirables-section-phantom]{Desiderata}
\item \hyperref[coding-section-phantom]{Style de programmation}
\item \hyperref[obsolete-section-phantom]{Obsolète}
\item \hyperref[fdl-section-phantom]{Licence de documentation libre GNU}
\item \hyperref[index-section-phantom]{Index généré automatiquement}
\end{enumerate}
\end{multicols}


\bibliography{my}
\bibliographystyle{amsalpha}

\end{document}
